% Section file for Chapter: Twisted Holography and Quantum Gravity
% Result (G3): Complementarity as Shifted-Symplectic Polarization

% Local macros (providecommand only; never \newcommand in chapter files)
\providecommand{\MC}{\mathrm{MC}}
\providecommand{\Defcyc}{\mathrm{Def}_{\mathrm{cyc}}}
\providecommand{\Definfmod}{\mathrm{Def}^{\mathrm{mod}}_\infty}
\providecommand{\Sh}{\mathrm{Sh}}
\providecommand{\gr}{\operatorname{gr}}
\providecommand{\id}{\mathrm{id}}
\providecommand{\Tr}{\operatorname{Tr}}
\providecommand{\Sym}{\operatorname{Sym}}
\providecommand{\Hom}{\operatorname{Hom}}
\providecommand{\End}{\operatorname{End}}
\providecommand{\Spec}{\operatorname{Spec}}
\providecommand{\Res}{\operatorname{Res}}
\providecommand{\rank}{\operatorname{rank}}
\providecommand{\ad}{\operatorname{ad}}
\providecommand{\Fred}{\operatorname{Fred}}
\providecommand{\Map}{\operatorname{Map}}
\providecommand{\RGamma}{R\Gamma}
\providecommand{\Perf}{\operatorname{Perf}}
\providecommand{\BTZ}{\mathrm{BTZ}}

\section{Complementarity as shifted-symplectic polarization}
% label removed: sec:thqg-symplectic-polarization
\label{ch:thqg-symplectic-polarization}% alias for dangling \ref{ch:...}
\index{symplectic polarization|textbf}
\index{complementarity!gravitational phase space}
\index{shifted symplectic!complementarity!holographic}

The complementarity theorem (Volume~I, Theorem~\ref{V1-thm:quantum-complementarity-main})
decomposes the ambient complex
$\mathbf{C}_g(\cA) = \RGamma(\overline{\mathcal{M}}_g,
\mathcal{Z}(\cA))$ into Lagrangian halves
$\mathbf{Q}_g(\cA) \oplus \mathbf{Q}_g(\cA^!)$.
Theorem~\ref{V1-thm:shifted-symplectic-complementarity} upgrades this to a
$(-1)$-shifted symplectic polarization via the BV formalism.
The ambient complex carries a Verdier involution whose eigenspace
decomposition is unconditional on the Koszul locus~(C1), while the
shifted-symplectic Lagrangian structure~(C2) is the natural
geometric home of the shadow obstruction tower
$\Theta_\cA^{\le 2} \to \Theta_\cA^{\le 3} \to \Theta_\cA^{\le 4}
\to \cdots$. We prove that the complementarity potential $S_\cA$
generates the dual Lagrangian as a formal graph, that its Taylor
jets are exactly the shadow jets
$H_\cA$, $\mathfrak{C}_\cA$, $\mathfrak{Q}_\cA$ of
Appendix~\ref*{V1-app:nonlinear-modular-shadows}, and that at genus~$1$
the construction reduces to a precise bulk--boundary entanglement
computation matching the BTZ partition function.

\medskip
\noindent\emph{Organization.}\;
\S\ref{subsec:thqg-III-ambient-complex} constructs the ambient
complex and Verdier involution, establishing the foundational
input.
\S\ref{subsec:thqg-III-eigenspace-decomposition} proves the
unconditional eigenspace decomposition~(C1) with complete detail.
\S\ref{subsec:thqg-III-shifted-symplectic} develops the
shifted-symplectic structure~(C2), including a self-contained
review of PTVV geometry.
\S\ref{subsec:thqg-III-complementarity-potential} constructs the
complementarity potential and identifies it with the shadow jet
expansion.
\S\ref{subsec:thqg-III-holographic-entanglement} gives the
holographic interpretation at genus~$1$ and the BTZ connection.
\S\ref{subsec:thqg-III-standard-landscape} verifies all structures
across the standard landscape.

% ======================================================================
%
% 1. THE AMBIENT COMPLEX AND VERDIER INVOLUTION
%
% ======================================================================

\subsection{The ambient complex and Verdier involution}
\label{subsec:thqg-III-ambient-complex}
\index{ambient complex|textbf}
\index{Verdier involution!on ambient complex}

The starting point is the datum of a chiral Koszul pair
$(\cA, \cA^!)$ on a smooth projective curve $X$ over $\mathbb{C}$.
The five main theorems (A--D+H) of Volume~I supply the following
chain of identifications.

\begin{definition}[Holographic ambient complex]
% label removed: def:thqg-III-holographic-ambient
\index{holographic ambient complex|textbf}
Let $(\cA, \cA^!)$ be a chiral Koszul pair and let $g \ge 0$.
The \emph{holographic ambient complex} at genus~$g$ is the cochain
complex
\begin{equation}% label removed: eq:thqg-III-holographic-ambient
\mathbf{C}_g(\cA)
:= \RGamma(\overline{\mathcal{M}}_g, \mathcal{Z}(\cA)),
\end{equation}
where $\mathcal{Z}(\cA)$ is the center local system on the
Deligne--Mumford--Knudsen compactification
$\overline{\mathcal{M}}_g$. When $g \ge 2$ and the relative bar
family is perfect
(Lemma~\ref{V1-lem:perfectness-criterion}), the fiber--center
identification
(Theorem~\ref{V1-thm:fiber-center-identification}) gives
\begin{equation}% label removed: eq:thqg-III-bar-center
\mathbf{C}_g(\cA)
\simeq
\RGamma(\overline{\mathcal{M}}_g,
R^0\pi_{g*}\barB^{(g)}(\cA)).
\end{equation}
\end{definition}

\begin{remark}[H/M/S layers]% label removed: rem:thqg-III-hms
Following Convention~\ref*{V1-conv:hms-levels}:
the complex $\mathbf{C}_g(\cA)$ is the H-level (homotopy)
object; its cohomology
$\mathcal{H}_g(\cA) := H^*(\mathbf{C}_g(\cA))
= H^*(\overline{\mathcal{M}}_g, \mathcal{Z}(\cA))$
is the S-level (shadow) object; and the bar complex
$(\barB^{(g)}(\cA), \Dg{g})$ is the M-level (model) object.
\end{remark}

\begin{proposition}[Properties of the holographic ambient complex;
\ClaimStatusConditional]
\label{prop:thqg-III-ambient-properties}
\index{holographic ambient complex!properties}
The holographic ambient complex
$\mathbf{C}_g(\cA)$ satisfies:
\begin{enumerate}[label=\textup{(\roman*)}]
\item \emph{Finite-dimensionality.}
 For each $n \in \mathbb{Z}$,
 $H^n(\mathbf{C}_g(\cA))$ is finite-dimensional over
 $\mathbb{C}$, and is nonzero only for
 $0 \le n \le 6g - 6$.

\item \emph{Verdier self-duality.}
 The Koszul pairing
 $\cA \otimes \cA^! \to \omega_X$ and the center isomorphism
 $Z(\cA) \cong Z(\cA^!)$
 (Lemma~\ref{V1-sublem:center-isomorphism}) induce a
 quasi-isomorphism of constructible sheaves on
 $\overline{\mathcal{M}}_g$:
 \begin{equation}% label removed: eq:thqg-III-verdier-self-duality
 \mathcal{Z}(\cA)
 \xrightarrow{\;\sim\;}
 \mathbb{D}\,\mathcal{Z}(\cA)[3g - 3],
 \end{equation}
 where $\mathbb{D}$ denotes Verdier duality on
 $\overline{\mathcal{M}}_g$.

\item \emph{Functoriality.}
 The assignment $(\cA, \cA^!) \mapsto \mathbf{C}_g(\cA)$
 is functorial in morphisms of Koszul pairs.
\end{enumerate}
\end{proposition}

\begin{proof}
\emph{Part (i).}
The cohomological dimension of $\overline{\mathcal{M}}_g$ is
$2 \dim_{\mathbb{C}} \overline{\mathcal{M}}_g = 2(3g - 3) = 6g - 6$
by Artin vanishing for proper DM stacks
\cite[\S4.1]{Olsson16}. The stalks of
$\mathcal{Z}(\cA)$ are finite-dimensional by hypothesis
(finite-dimensional fiber cohomology,
Lemma~\ref{V1-lem:perfectness-criterion}(ii)).
By constructibility of $\mathcal{Z}(\cA)$
(Lemma~\ref{V1-lem:quantum-ss-convergence}, condition (2)),
each cohomology sheaf $R^q\pi_{g*}\barB^{(g)}(\cA)$
is a constructible sheaf on the Noetherian stack
$\overline{\mathcal{M}}_g$. Cohomology of a constructible
sheaf on a Noetherian DM stack is finite-dimensional
\cite[Theorem~18.1.1]{SAG}.

\emph{Part (ii).}
The Koszul pairing restricts to a perfect pairing on centers
$Z(\cA) \otimes Z(\cA^!) \to \mathbb{C}$
(Corollary~\ref{V1-cor:duality-bar-complexes-complete} and
Lemma~\ref{V1-sublem:center-isomorphism}).
The center isomorphism $Z(\cA) \cong Z(\cA^!)$ converts the
pairing into a self-pairing on $\mathcal{Z}(\cA)$.
By Verdier duality for constructible sheaves on smooth proper
DM stacks, a nondegenerate self-pairing on $\mathcal{Z}(\cA)$
is equivalent to an isomorphism
$\mathcal{Z}(\cA) \xrightarrow{\sim}
\mathbb{D}\,\mathcal{Z}(\cA)[3g - 3]$.

\emph{Part (iii).}
The center local system $\mathcal{Z}(-)$ is functorial in
morphisms of chiral algebras (being a homotopy fiber of the
bar complex, which is functorial by
Theorem~\ref{V1-thm:geometric-equals-operadic-bar}).
Applying $\RGamma(\overline{\mathcal{M}}_g, -)$ preserves
functoriality.
\end{proof}

\begin{remark}[Attribution: Vol~I dependencies for
Proposition~\ref{prop:thqg-III-ambient-properties}]
\label{rem:thqg-III-ambient-properties-attribution}
The proof of
Proposition~\ref{prop:thqg-III-ambient-properties} is
\emph{conditional} on the following Volume~I ingredients, cited at
the citation level rather than re-inscribed here: the perfectness
criterion for the relative bar family
(Lemma~\ref{V1-lem:perfectness-criterion}), the quantum
semisimplicity convergence lemma supplying constructibility of the
center local system
(Lemma~\ref{V1-lem:quantum-ss-convergence}), the center--center
isomorphism
(Lemma~\ref{V1-sublem:center-isomorphism}), the duality of
complete bar complexes
(Corollary~\ref{V1-cor:duality-bar-complexes-complete}), and the
geometric--operadic bar comparison
(Theorem~\ref{V1-thm:geometric-equals-operadic-bar}). Each of
these is a Volume~I result; this proposition inherits their
scopes (Koszul locus, finite-dimensional fiber cohomology). The
$\ClaimStatusConditional$ tag records this external dependency
per HZ-11 and AP287.
\end{remark}

\begin{construction}[Verdier involution on the ambient complex]
\label{constr:thqg-III-verdier-involution}
\index{Verdier involution!construction|textbf}
The Verdier self-duality of
Proposition~\ref{prop:thqg-III-ambient-properties}(ii) and
the Koszul involutivity
$(\cA^!)^! \simeq \cA$
(Theorem~\ref{V1-thm:chiral-koszul-duality}) combine to produce
a cochain-level endomorphism
\begin{equation}% label removed: eq:thqg-III-sigma-def
\sigma
\colon
\mathbf{C}_g(\cA) \longrightarrow \mathbf{C}_g(\cA)
\end{equation}
constructed as the composition:
\begin{equation}% label removed: eq:thqg-III-sigma-composition
\begin{tikzcd}[column sep=large]
\mathbf{C}_g(\cA)
\arrow[r, "\mathbb{D}"]
&
\mathbf{C}_g(\cA^!)^{\vee}[{-(3g{-}3)}]
\arrow[r, "\mathrm{KS}^{\vee}"]
&
\mathbf{C}_g(\cA)^{\vee\vee}[{-(3g{-}3)}][{3g{-}3}]
\arrow[r, "\mathrm{ev}"]
&
\mathbf{C}_g(\cA)
\end{tikzcd}
\end{equation}
Here $\mathbb{D}$ is the Verdier duality map from
Corollary~\ref{V1-cor:duality-bar-complexes-complete},
$\mathrm{KS}$ is the Koszul identification
$\mathbf{C}_g(\cA) \xrightarrow{\sim} \mathbf{C}_g(\cA^!)$
via Lemma~\ref{V1-sublem:center-isomorphism}, and $\mathrm{ev}$
is the evaluation map. The shifts cancel, and the composition
is well-defined as a cochain map of degree zero.
\end{construction}

\begin{proposition}[Involutivity and anti-symmetry;
\ClaimStatusProvedHere]
\label{prop:thqg-III-involutivity}
\index{Verdier involution!involutivity}
The endomorphism $\sigma$ of
Construction~\textup{\ref{constr:thqg-III-verdier-involution}}
satisfies:
\begin{enumerate}[label=\textup{(\alph*)}]
\item $\sigma^2 = \id$ \textup{(}involutivity\textup{)}.
\item The Verdier pairing
 $\langle -, - \rangle_{\mathbb{D}} \colon
 \mathbf{C}_g(\cA) \otimes \mathbf{C}_g(\cA) \to
 \mathbb{C}[{-(3g{-}3)}]$
 satisfies the anti-symmetry
 \begin{equation}% label removed: eq:thqg-III-anti-symmetry
 \langle \sigma(v), \sigma(w) \rangle_{\mathbb{D}}
 = -\langle v, w \rangle_{\mathbb{D}}
 \end{equation}
 for all $v, w \in \mathbf{C}_g(\cA)$.
\end{enumerate}
\end{proposition}

\begin{proof}
\emph{Part (a).}
Koszul involutivity
$(\cA^!)^! \simeq \cA$ gives $\mathrm{KS}^{-1} = \mathrm{KS}'$
(the transpose identification). Verdier duality satisfies
$\mathbb{D}^2 = \id$ on constructible sheaves over
smooth proper DM stacks (Poincar\'{e} duality is involutive).
The composition $\sigma$ is therefore the identity composed
twice, once from each factor:
\[
\sigma^2
= (\mathrm{ev} \circ \mathrm{KS}'^{\vee} \circ \mathbb{D})
 \circ
 (\mathrm{ev} \circ \mathrm{KS}^{\vee} \circ \mathbb{D})
= \id.
\]

\emph{Part (b).}
Let $v, w \in \mathbf{C}_g(\cA)$. By definition of the
Verdier pairing:
\[
\langle \sigma(v), \sigma(w) \rangle_{\mathbb{D}}
= \langle \mathbb{D}(v'), \mathbb{D}(w') \rangle_{\mathbb{D}}
\]
where $v' = \mathrm{KS}(v)$ and $w' = \mathrm{KS}(w)$.
The Verdier pairing satisfies $\mathbb{D}$ anti-commutativity
with the Kodaira--Spencer action
(Theorem~\ref{V1-thm:kodaira-spencer-chiral-complete},
equation~\eqref{V1-eq:verdier-ks-anticommute}). At the level of
the bilinear form, this gives
$\langle \mathbb{D}(x), \mathbb{D}(y) \rangle_{\mathbb{D}}
= -\langle x, y \rangle_{\mathbb{D}}$
(Proposition~\ref{V1-prop:lagrangian-eigenspaces}(ii)).
Composing with the Koszul identification (which preserves the
pairing up to sign by construction) gives the stated
anti-symmetry.
\end{proof}

\begin{proposition}[Compatibility with the modular MC element;
\ClaimStatusProvedHere]
\label{prop:thqg-III-mc-compatibility}
\index{Maurer--Cartan element!Verdier compatibility}
Let $\Theta_\cA \in \MC(\mathfrak{g}_\cA^{\mathrm{mod}})$ be the
bar-intrinsic MC element
\textup{(}Theorem~\textup{\ref*{V1-thm:mc2-bar-intrinsic})}.
The Verdier involution $\sigma$ intertwines the shadow projections
of $\Theta_\cA$ and $\Theta_{\cA^!}$:
\begin{equation}% label removed: eq:thqg-III-theta-intertwine
\sigma \circ \pi_g^*(\Theta_\cA)
= -\pi_g^*(\Theta_{\cA^!}),
\end{equation}
where $\pi_g^*$ denotes the genus-$g$ projection.
\end{proposition}

\begin{proof}
The MC element $\Theta_\cA := D_\cA - d_0$ satisfies the MC
equation because $D_\cA^2 = 0$
(Theorem~\ref*{V1-thm:mc2-bar-intrinsic}).
Under Verdier--Koszul duality
(Theorem~\ref*{V1-thm:verdier-bar-cobar}), the bar
coderivation $D_\cA$ maps to $D_{\cA^!}$ while the genus-$0$
collision differential $d_0$ is invariant. At the level of
coderivations, therefore,
$\sigma(\Theta_\cA) = D_{\cA^!} - d_0 = \Theta_{\cA^!}$.
The sign in~\eqref{V1-eq:thqg-III-theta-intertwine} arises when
one passes to genus-$g$ projections through the bilinear
pairing. By
Proposition~\ref{V1-prop:lagrangian-eigenspaces}(ii), $\sigma$
acts anti-symmetrically on the pairing:
$\langle \sigma(x), \sigma(y) \rangle
= -\langle x, y \rangle$.
The projection $\pi_g^*$ extracts the component in Hodge
degree~$g$ via this pairing, so the anti-symmetry contributes
an overall sign:
\[
\sigma \circ \pi_g^*(\Theta_\cA)
= -\pi_g^*(D_{\cA^!} - d_0)
= -\pi_g^*(\Theta_{\cA^!}).
\]
This is the complementarity relation
$\pi_g^*(\Theta_\cA) + \pi_g^*(\Theta_{\cA^!}) = 0$
at every genus.
\end{proof}

\begin{remark}[The chain-level mechanism]
% label removed: rem:thqg-III-chain-level
\index{complementarity!chain-level mechanism}
Proposition~\ref{prop:thqg-III-mc-compatibility} is the
chain-level incarnation of the complementarity relation
$\sigma \circ \pi_g^*(\Theta_\cA) = -\pi_g^*(\Theta_{\cA^!})$
(Remark~\ref{V1-rem:theorem-C-decomposition}).
At the S-level, the pairing-projected relation becomes
$\kappa(\cA) + \kappa(\cA^!) = 0$ for Kac--Moody and free-field algebras,
but $\kappa(\cA) + \kappa(\cA^!) = 13$ for Virasoro and
$\kappa + \kappa^! = \varrho_{\cA}\,K$ for general $\mathcal{W}$-algebras
(with $\varrho_{\cW_N} = H_N - 1$ and $K = 4N^3 - 2N - 2$,
$\varrho_{\BP} = 1/6$ and $K_{\BP} = 196$):
the complementarity sum is family-dependent, not universally
zero (Volume~I, Theorem~D).
At the H-level, $\sigma$ intertwines the two MC elements in
opposite sign, implementing the deformation--obstruction
exchange.
\end{remark}

\begin{definition}[Holographic Verdier pairing]
\label{def:thqg-III-holographic-pairing}
\index{Verdier pairing!holographic|textbf}
The \emph{holographic Verdier pairing} at genus~$g$ is the
bilinear form
\begin{equation}% label removed: eq:thqg-III-holographic-pairing-def
\langle -, - \rangle_g
\colon
\mathbf{C}_g(\cA) \otimes \mathbf{C}_g(\cA)
\longrightarrow
\mathbb{C}[{-(3g{-}3)}]
\end{equation}
induced by the Verdier self-duality
\eqref{V1-eq:thqg-III-verdier-self-duality} and the $\RGamma$
functor. Explicitly, for cochains
$\alpha \in C^p(\overline{\mathcal{M}}_g, \mathcal{Z}(\cA))$
and
$\beta \in C^{(3g-3)-p}(\overline{\mathcal{M}}_g,
\mathcal{Z}(\cA))$,
\begin{equation}% label removed: eq:thqg-III-pairing-explicit
\langle \alpha, \beta \rangle_g
:=
\int_{\overline{\mathcal{M}}_g}
\alpha \cup \beta
\;\in\; \mathbb{C},
\end{equation}
where the cup product uses the self-pairing on
$\mathcal{Z}(\cA)$ and integration uses the fundamental class
of $\overline{\mathcal{M}}_g$.
\end{definition}

\begin{lemma}[Nondegeneracy of the holographic pairing;
\ClaimStatusProvedHere]
\label{lem:thqg-III-nondegeneracy}
\index{Verdier pairing!nondegeneracy}
When the relative bar family is perfect
\textup{(}Lemma~\textup{\ref{V1-lem:perfectness-criterion})},
the holographic Verdier pairing
\eqref{V1-eq:thqg-III-holographic-pairing-def} is nondegenerate.
\end{lemma}

\begin{proof}
Perfectness of $R\pi_{g*}\barB^{(g)}(\cA)$ ensures that
the Verdier duality map
$\mathcal{Z}(\cA) \xrightarrow{\sim}
\mathbb{D}\,\mathcal{Z}(\cA)[3g - 3]$
is a quasi-isomorphism of perfect complexes on the smooth
proper DM stack $\overline{\mathcal{M}}_g$. Applying $\RGamma$
and Serre duality on the stack gives a perfect pairing on
cohomology. On the cochain level, the cup-product pairing
\[
C^p(\overline{\mathcal{M}}_g, \mathcal{Z})
\otimes
C^{(3g-3)-p}(\overline{\mathcal{M}}_g, \mathbb{D}\mathcal{Z})
\to
C^{3g-3}(\overline{\mathcal{M}}_g, \omega_{\overline{\mathcal{M}}_g})
\xrightarrow{\int}
\mathbb{C}
\]
is nondegenerate by the duality theorem for constructible sheaves
on smooth proper varieties
(Grothendieck--Verdier duality, \cite[Theorem~3.1.10]{KS90}).
\end{proof}

\begin{proposition}[Degree shift at each genus;
\ClaimStatusProvedHere]
% label removed: prop:thqg-III-degree-shift
The holographic Verdier pairing has cohomological degree
$-(3g - 3)$. In particular:
\begin{enumerate}[label=\textup{(\alph*)}]
\item At genus $0$: degree $+3$, pairing
 $H^0 \otimes H^0 \to \mathbb{C}[3]$
 \textup{(}trivially degenerate on one-dimensional $H^0$;
 the complementarity is purely at the S-level\textup{)}.
\item At genus $1$: degree $0$, pairing
 $H^0 \otimes H^0 \to \mathbb{C}$
 \textup{(}classical pairing\textup{)}.
\item At genus $g \ge 2$: degree $-(3g - 3)$,
 giving a genuine shifted-symplectic structure.
\end{enumerate}
\end{proposition}

\begin{proof}
The degree is $-(3g - 3) = -\dim_{\mathbb{C}}
\overline{\mathcal{M}}_g$, which is the Verdier shift for
Serre duality on a smooth variety of dimension $3g - 3$.
The three cases follow by substituting $g = 0, 1, \ge 2$.
\end{proof}

\begin{remark}[Genus $1$ distinguished]
% label removed: rem:thqg-III-genus-1-special
\index{genus 1!holographic pairing}
At genus~$1$, $\dim \overline{\mathcal{M}}_{1,1} = 1$
but $\dim \overline{\mathcal{M}}_1 = 0$ (a point for the
coarse moduli). The relevant moduli space for marked curves
is $\overline{\mathcal{M}}_{1,1}$, where the pairing has
degree~$0$. This is the reason genus~$1$ complementarity
has the simplest form:
$Q_1(\cA) \oplus Q_1(\cA^!) \cong H^*(\overline{\mathcal{M}}_{1,1},
\mathcal{Z}(\cA))$ with a classical (unshifted) pairing.
\end{remark}

% ======================================================================
%
% 2. THE UNCONDITIONAL EIGENSPACE DECOMPOSITION (C1)
%
% ======================================================================

\subsection{The unconditional eigenspace decomposition}
\label{subsec:thqg-III-eigenspace-decomposition}
\index{eigenspace decomposition!holographic|textbf}
\index{complementarity!C1 decomposition}

The eigenspace decomposition~(C1) holds unconditionally on the
Koszul locus, requiring no perfectness or nondegeneracy hypotheses.
We give a self-contained proof in the holographic setting.

\begin{theorem}[Holographic eigenspace decomposition (C1);
\ClaimStatusProvedHere]
\label{thm:thqg-III-eigenspace-decomposition}
\index{Lagrangian!eigenspace decomposition!holographic}
Let $(\cA, \cA^!)$ be a chiral Koszul pair on a smooth projective
curve $X$ over $\mathbb{C}$, and let $g \ge 0$.

\smallskip\noindent
\textbf{H-level.}\;
The Verdier involution $\sigma$ of
Construction~\textup{\ref{constr:thqg-III-verdier-involution}}
splits the ambient complex:
\begin{equation}% label removed: eq:thqg-III-C1-homotopy
\mathbf{C}_g(\cA)
\;\simeq\;
\mathbf{Q}_g(\cA) \;\oplus\; \mathbf{Q}_g(\cA^!),
\end{equation}
where
$\mathbf{Q}_g(\cA) := \operatorname{fib}(\sigma - \id)$ and
$\mathbf{Q}_g(\cA^!) := \operatorname{fib}(\sigma + \id)$
are the homotopy eigenspaces.

\smallskip\noindent
\textbf{S-level.}\;
On cohomology:
\begin{equation}% label removed: eq:thqg-III-C1-shadow
\mathcal{H}_g(\cA)
:= H^*(\overline{\mathcal{M}}_g, \mathcal{Z}(\cA))
\;=\;
Q_g(\cA) \;\oplus\; Q_g(\cA^!),
\end{equation}
where $Q_g(\cA) = \ker(\sigma - \id)$ and
$Q_g(\cA^!) = \ker(\sigma + \id)$.

\smallskip\noindent
\textbf{Duality.}\;
The Verdier pairing restricts to a perfect duality
\begin{equation}% label removed: eq:thqg-III-C1-duality
Q_g(\cA) \;\cong\; Q_g(\cA^!)^{\vee}.
\end{equation}

\smallskip\noindent
\textbf{Functoriality.}\;
The decomposition is natural in morphisms of Koszul pairs and
compatible with the conformal weight and cohomological degree
gradings.
\end{theorem}

\begin{proof}
The proof proceeds in four stages.

\emph{Stage 1: Projector construction.}
Since $\sigma^2 = \id$
(Proposition~\ref{prop:thqg-III-involutivity}(a)) and we
work over $\mathbb{C}$ (characteristic $\ne 2$), the cochain
maps
\begin{equation}% label removed: eq:thqg-III-projectors
p^+ := \tfrac{1}{2}(\id + \sigma),
\qquad
p^- := \tfrac{1}{2}(\id - \sigma)
\end{equation}
are idempotent cochain endomorphisms of $\mathbf{C}_g(\cA)$
satisfying:
\begin{align}
(p^+)^2 &= p^+, &
(p^-)^2 &= p^-, &
p^+ + p^- &= \id, &
p^+ \circ p^- &= 0.
% label removed: eq:thqg-III-projector-properties
\end{align}
These are immediate from $\sigma^2 = \id$. Explicitly:
\[
(p^+)^2
= \tfrac{1}{4}(\id + \sigma)^2
= \tfrac{1}{4}(\id + 2\sigma + \sigma^2)
= \tfrac{1}{4}(\id + 2\sigma + \id)
= \tfrac{1}{2}(\id + \sigma)
= p^+.
\]
The verification for $p^-$ is identical. For the product:
\[
p^+ \circ p^-
= \tfrac{1}{4}(\id + \sigma)(\id - \sigma)
= \tfrac{1}{4}(\id - \sigma^2)
= 0.
\]
Since $\sigma$ commutes with the differential
$d$ on $\mathbf{C}_g(\cA)$ (it is a cochain map by
construction), both $p^+$ and $p^-$ are cochain maps.

\emph{Stage 2: Direct sum decomposition.}
Define
\[
\mathbf{C}_g^+ := \operatorname{im}(p^+),
\qquad
\mathbf{C}_g^- := \operatorname{im}(p^-).
\]
The idempotent decomposition $p^+ + p^- = \id$ gives
\begin{equation}% label removed: eq:thqg-III-direct-sum
\mathbf{C}_g(\cA) = \mathbf{C}_g^+ \oplus \mathbf{C}_g^-
\end{equation}
as cochain complexes (not merely as graded vector spaces,
because $p^\pm$ are cochain maps). Since
$\sigma|_{\mathbf{C}_g^+} = +\id$ and
$\sigma|_{\mathbf{C}_g^-} = -\id$, the subcomplexes
$\mathbf{C}_g^+$ and $\mathbf{C}_g^-$ are the $+1$ and $-1$
eigenspaces of $\sigma$.

It remains to identify $\mathbf{C}_g^+$ with the homotopy
eigenspace $\mathbf{Q}_g(\cA) = \operatorname{fib}(\sigma - \id)$.
The cone sequence
$\operatorname{fib}(\sigma - \id) \to \mathbf{C}_g
\xrightarrow{\sigma - \id} \mathbf{C}_g$
gives $\operatorname{fib}(\sigma - \id) \simeq
\ker(\sigma - \id) = \operatorname{im}(p^+) = \mathbf{C}_g^+$,
where the second identification uses that $\sigma - \id$ is
split-surjective onto $\mathbf{C}_g^-$ (the splitting is given
by $-\tfrac{1}{2}\id$ restricted to $\mathbf{C}_g^-$).
Similarly $\operatorname{fib}(\sigma + \id) \simeq \mathbf{C}_g^-$.

\emph{Stage 3: Eigenspace identification.}
We identify $\mathbf{C}_g^+ \simeq \mathbf{Q}_g(\cA)$ and
$\mathbf{C}_g^- \simeq \mathbf{Q}_g(\cA^!)$.
The $+1$-eigenspace $\mathbf{C}_g^+$ consists of cochains
$\alpha$ with $\sigma(\alpha) = \alpha$. By construction
of $\sigma$
(Construction~\ref{constr:thqg-III-verdier-involution}),
$\sigma(\alpha) = \alpha$ means that $\alpha$ is fixed under
the Verdier--Koszul duality, i.e., $\alpha$ comes from the
$j_*$-extension of bar cochains of $\cA$
(Lemma~\ref{V1-lem:eigenspace-decomposition-complete}).
The genus-$g$ quantum correction space $Q_g(\cA)$ is precisely
the subspace of $H^*(\overline{\mathcal{M}}_g, \mathcal{Z}(\cA))$
arising from bar cochains of $\cA$ in the $j_*$-extension
(Lemma~\ref{V1-lem:quantum-from-ss}).

Symmetrically, $\mathbf{C}_g^-$ consists of cochains fixed by
$-\sigma$, which come from the $j_!$-extension of bar
cochains of $\cA^!$ (equivalently, cobar cochains of $\cA$),
giving $\mathbf{C}_g^- \simeq \mathbf{Q}_g(\cA^!)$.

\emph{Stage 4: S-level decomposition and duality.}
Applying $H^*$ to the direct sum
\eqref{V1-eq:thqg-III-direct-sum} gives:
\[
H^*(\mathbf{C}_g(\cA))
= H^*(\mathbf{C}_g^+) \oplus H^*(\mathbf{C}_g^-)
= Q_g(\cA) \oplus Q_g(\cA^!).
\]
The duality $Q_g(\cA) \cong Q_g(\cA^!)^{\vee}$ follows from
the Verdier pairing. For $v \in Q_g(\cA)$ and
$w \in Q_g(\cA^!)$, the pairing $\langle v, w \rangle_g$
is nondegenerate: if $\langle v, w \rangle_g = 0$ for all
$w \in Q_g(\cA^!)$, then $v$ lies in the radical of the
restricted pairing. Since
$\langle Q_g(\cA), Q_g(\cA) \rangle_g = 0$
(isotropy from anti-symmetry of $\sigma$),
$v$ lies in the radical of the full pairing on
$\mathcal{H}_g(\cA)$, which is zero by nondegeneracy
(Lemma~\ref{lem:thqg-III-nondegeneracy}).

\emph{Note on unconditionality.}
Stages 1--3 use only $\sigma^2 = \id$ and the
identification of $\sigma$-eigenspaces with $j_*/j_!$
extensions. No perfectness or nondegeneracy hypothesis is
needed. The duality statement in Stage~4 uses
nondegeneracy but can also be obtained unconditionally from
spectral sequence duality
(Corollary~\ref{V1-cor:quantum-dual-complete}).

\emph{Functoriality.}
A morphism $f \colon (\cA, \cA^!) \to (\cB, \cB^!)$
of Koszul pairs induces a map
$f_* \colon \mathbf{C}_g(\cA) \to \mathbf{C}_g(\cB)$
by functoriality of the center local system and $\RGamma$.
Since $f$ commutes with Koszul duality
($f^! \colon \cA^! \to \cB^!$ is the Koszul dual map),
$f_*$ commutes with $\sigma$, hence preserves the
projectors $p^\pm$ and the eigenspace decomposition.
\end{proof}

\begin{corollary}[Complementarity exchange principle;
\ClaimStatusProvedHere]
% label removed: cor:thqg-III-complementarity-exchange
\index{complementarity!exchange principle}
For any functional $\Phi$ on chiral Koszul pairs that
factors through the ambient complex, the decomposition
\[
\Phi(\cA) + \Phi(\cA^!) = \Phi_{\mathrm{total}}
\]
holds with $\Phi_{\mathrm{total}}$ independent of the choice
of $\cA$ within its Koszul pair.
\end{corollary}

\begin{proof}
Since $\Phi$ factors through
$\mathcal{H}_g(\cA) = Q_g(\cA) \oplus Q_g(\cA^!)$,
the restriction of $\Phi$ to each summand gives the
decomposition. Independence of $\Phi_{\mathrm{total}}$
from the choice of $\cA$ within $\{\cA, \cA^!\}$
follows from the symmetry $Q_g(\cA) \oplus Q_g(\cA^!)
= Q_g(\cA^!) \oplus Q_g((\cA^!)^!) = Q_g(\cA^!) \oplus Q_g(\cA)$.
\end{proof}

\begin{corollary}[Dimension parity;
\ClaimStatusProvedHere]
% label removed: cor:thqg-III-dimension-parity
For a Koszul self-dual chiral algebra $\cA \cong \cA^!$ and
$g \ge 1$, the total dimension
$\dim \mathcal{H}_g(\cA)$ is even.
\end{corollary}

\begin{proof}
Self-duality gives $Q_g(\cA) \cong Q_g(\cA^!)$, so
$\dim Q_g(\cA) = \dim Q_g(\cA^!)$ and
$\dim \mathcal{H}_g(\cA)
= 2 \dim Q_g(\cA)$
(Theorem~\ref{thm:self-dual-halving}).
\end{proof}

\begin{proposition}[Genus-$0$ complementarity;
\ClaimStatusProvedHere]
% label removed: prop:thqg-III-genus-0
\index{complementarity!genus 0}
At genus~$0$, $\overline{\mathcal{M}}_0 = \mathrm{pt}$ and
$\mathcal{H}_0(\cA) = Z(\cA)$.
The Verdier involution $\sigma$ acts as $+\id$ on $Z(\cA)$,
giving:
\[
Q_0(\cA) = Z(\cA), \qquad Q_0(\cA^!) = 0.
\]
\end{proposition}

\begin{proof}
Since $\overline{\mathcal{M}}_0$ is a point, the center
local system has a single stalk $Z(\cA)$. The bar complex
at genus~$0$ is uncurved (the curvature
$\dfib^{\,2} = \kappa(\cA) \cdot \omega_g$ vanishes
because $\omega_0 = 0$), so the Verdier involution acts on
the genus-$0$ center by the identity: bar cochains with
$j_*$-extension are all of $Z(\cA)$, and there are no
$j_!$-extended cochains (the configuration space has no
boundary at genus~$0$).
\end{proof}

\begin{proposition}[Genus-$1$ complementarity;
\ClaimStatusProvedHere]
\label{prop:thqg-III-genus-1}
\index{complementarity!genus 1}
At genus~$1$,
$H^*(\overline{\mathcal{M}}_{1,1})
= \mathbb{Q} \oplus \mathbb{Q}\lambda_1$
with $\lambda_1 = c_1(\mathbb{E})$ the Hodge class.
For the Heisenberg algebra $\mathcal{H}_\kappa$:
\begin{equation}% label removed: eq:thqg-III-heisenberg-genus-1
Q_1(\mathcal{H}_\kappa) = \mathbb{C} \cdot \kappa,
\qquad
Q_1(\mathcal{H}_\kappa^!) = \mathbb{C} \cdot \lambda_1.
\end{equation}
The central extension $\kappa$ is an obstruction for
$\mathcal{H}_\kappa$ (eigenvalue $+1$); the Hodge class
$\lambda_1$ is the curvature for $\mathcal{H}_\kappa^!$
(eigenvalue $-1$).
\end{proposition}

\begin{proof}
The involution $\sigma$ acts on the
two-dimensional space $H^0 \oplus H^2$ of
$\overline{\mathcal{M}}_{1,1}$. The class $\kappa$ is the
image of the level parameter under the genus-$1$ bar
complex. Since $\kappa$
arises from $j_*$-extended bar cochains, it has eigenvalue
$+1$. The class $\lambda_1$ is the first Chern class of
the Hodge bundle, which arises from the curvature of the
genus-$1$ Gauss--Manin connection; it comes from $j_!$-extended
cochains and has eigenvalue $-1$.

Dimension check: $\dim Q_1 + \dim Q_1^! = 1 + 1 = 2
= \dim H^*(\overline{\mathcal{M}}_{1,1})$.
\end{proof}

\begin{remark}[Deformation--obstruction exchange at genus $1$]
% label removed: rem:thqg-III-def-obs-exchange
\index{deformation--obstruction exchange!genus 1}
The genus-$1$ complementarity is the simplest nontrivial
instance of the exchange principle:
\begin{center}
\begin{tabular}{l@{\quad}c@{\quad}c}
& $\mathcal{H}_\kappa$ & $\mathcal{H}_\kappa^!$ \\
\midrule
Obstruction & $\kappa$ & --- \\
Deformation & --- & $\lambda_1$ \\
\end{tabular}
\end{center}
What $\mathcal{H}_\kappa$ sees as the obstruction to
extending its genus-$0$ structure to genus~$1$ (the
central extension $\kappa$), the Koszul dual
$\mathcal{H}_\kappa^!$ sees as a deformation (the Hodge
class $\lambda_1$ parametrizing the curvature of the
genus-$1$ structure).
\end{remark}

% ======================================================================
%
% 3. SHIFTED-SYMPLECTIC STRUCTURE (C2)
%
% ======================================================================

\subsection{The shifted-symplectic structure}
\label{subsec:thqg-III-shifted-symplectic}
\index{shifted symplectic!structure!holographic|textbf}
\index{PTVV!holographic application}

The unconditional decomposition~(C1) is an eigenspace statement.
The conditional upgrade~(C2) endows the ambient complex with a
$(-1)$-shifted symplectic structure in the sense of
Pantev--To\"{e}n--Vaqui\'{e}--Vezzosi \cite{PTVV13}, making
$Q_g(\cA)$ and $Q_g(\cA^!)$ complementary Lagrangians in the
precise derived-algebraic sense.

\subsubsection{Review of PTVV shifted symplectic geometry}
% label removed: subsubsec:thqg-III-ptvv-review
\index{PTVV!shifted symplectic geometry}

We recall the essential definitions from \cite{PTVV13} in the
form needed for our application.

\begin{definition}[{$n$-shifted symplectic structure
\cite[Definition~1.18]{PTVV13}}]
\label{def:thqg-III-shifted-symplectic}
\index{shifted symplectic!structure|textbf}
Let $\mathbf{F}$ be a derived stack locally of finite
presentation. An \emph{$n$-shifted symplectic structure} on
$\mathbf{F}$ is a closed $2$-form
$\omega \in \mathcal{A}^{2,\mathrm{cl}}(\mathbf{F}, n)$
of degree $n$ such that the induced map
\begin{equation}% label removed: eq:thqg-III-ptvv-nondegen
\omega^{\sharp}
\colon
\mathbb{T}_{\mathbf{F}} \longrightarrow
\mathbb{L}_{\mathbf{F}}[n]
\end{equation}
is a quasi-isomorphism, where $\mathbb{T}_{\mathbf{F}}$ and
$\mathbb{L}_{\mathbf{F}}$ are the tangent and cotangent
complexes.
\end{definition}

\begin{definition}[{Lagrangian in the PTVV sense
\cite[Definition~2.8]{PTVV13}}]
\label{def:thqg-III-ptvv-lagrangian}
\index{Lagrangian!PTVV sense|textbf}
Let $(\mathbf{F}, \omega)$ be an $n$-shifted symplectic derived
stack. A morphism $f \colon \mathbf{L} \to \mathbf{F}$ is
\emph{Lagrangian} if:
\begin{enumerate}[label=\textup{(\alph*)}]
\item \emph{Isotropy:} $f^* \omega = 0$ in
 $\mathcal{A}^{2,\mathrm{cl}}(\mathbf{L}, n)$.
\item \emph{Nondegeneracy:} the induced map
 $\mathbb{T}_{\mathbf{L}} \to
 \mathbb{L}_{\mathbf{L}/\mathbf{F}}[n-1]$
 is a quasi-isomorphism.
\end{enumerate}
Equivalently, condition (b) says the induced map
$\mathbf{L} \to (\mathbf{F}/\mathbf{L})^{\vee}[n]$ is a
quasi-isomorphism, where $\mathbf{F}/\mathbf{L}$ is the
cofiber.
\end{definition}

\begin{example}[Linear shifted symplectic space]
\label{ex:thqg-III-linear-shifted}
\index{shifted symplectic!linear example}
For a cochain complex $V$ with a nondegenerate symmetric pairing
$\langle -, - \rangle \colon V \otimes V \to \mathbb{C}[n]$,
the associated derived affine scheme
$\mathbf{V} = \Spec(\Sym(V^*))$ carries an $n$-shifted
symplectic structure: the pairing defines a constant $2$-form
of degree $n$, and the closure condition
$d_{\mathrm{dR}}\omega = 0$ is automatic for constant forms
on linear spaces \cite[Example~1.4]{PTVV13}.

A subcomplex $L \hookrightarrow V$ is Lagrangian if and only
if:
\begin{enumerate}[label=\textup{(\alph*)}]
\item $\langle L, L \rangle = 0$ (isotropy), and
\item $L \to (V/L)^{\vee}[n]$ is a quasi-isomorphism
 (nondegeneracy).
\end{enumerate}
\end{example}

\begin{proposition}[{Kontsevich--Pridham correspondence
\cite{Pridham17}}]
\label{prop:thqg-III-kontsevich-pridham}
\index{Kontsevich--Pridham correspondence}
Let $\mathfrak{g}$ be a dg Lie algebra with a nondegenerate
invariant pairing of degree $n$:
\[
\langle -, - \rangle
\colon \mathfrak{g} \otimes \mathfrak{g}
\to \mathbb{C}[n],
\qquad
\langle [x,y], z \rangle
= \langle x, [y,z] \rangle.
\]
Then the formal moduli problem
$\MC(\mathfrak{g})$ carries an $n$-shifted symplectic structure.
\end{proposition}

\begin{remark}[Two shifted structures in complementarity]
% label removed: rem:thqg-III-two-structures
\index{shifted symplectic!two structures}
Complementarity carries two shifted symplectic structures:
\begin{enumerate}[label=\textup{(\roman*)}]
\item The \emph{BV structure}: the antibracket on
 $\barB^{\mathrm{ch}}(\cA)$ gives a $(-1)$-shifted symplectic
 structure on $\MC(\barB^{(g)}(\cA)[1])$, universal in $g$.
\item The \emph{Verdier structure}: the holographic pairing on
 $\mathbf{C}_g(\cA)$ gives a $(-(3g-3))$-shifted symplectic
 structure, depending on $g$.
\end{enumerate}
Both make $Q_g(\cA)$ and $Q_g(\cA^!)$ complementary Lagrangians;
the two are compatible via the bar-complex spectral sequence.
\end{remark}

\subsubsection{The BV shifted-symplectic structure}
% label removed: subsubsec:thqg-III-bv-shifted
\index{BV algebra!shifted symplectic|textbf}

\begin{theorem}[BV Lagrangian polarization;
\ClaimStatusProvedHere]
\label{thm:thqg-III-bv-lagrangian}
\index{BV algebra!Lagrangian polarization}
Let $(\cA, \cA^!)$ be a chiral Koszul pair and $g \ge 1$.
\begin{enumerate}[label=\textup{(\roman*)}]
\item The dg Lie algebra
 $L_g := \barB^{(g)}(\cA)[1]$ carries a nondegenerate
 invariant pairing of degree $-1$, making $\MC(L_g)$
 a $(-1)$-shifted symplectic formal moduli problem.

\item The $\sigma$-eigenspace decomposition
 $L_g = L_g^+ \oplus L_g^-$ provides complementary
 Lagrangian subspaces.

\item $L_g^+$ controls deformations of $\cA$ at genus $g$;
 $L_g^-$ controls deformations of $\cA^!$ at genus $g$.
\end{enumerate}
\end{theorem}

\begin{proof}
\emph{Part (i).}
The BV antibracket $\{-, -\}_{\mathrm{BV}}$ on
$\barB^{\mathrm{ch}}(\cA)$ has degree $+1$
(Theorem~\ref{V1-thm:config-space-bv}). Shifting by $[1]$
converts this to a degree-$0$ Lie bracket on $L_g$.
The BV pairing on $\barB^{(g)}(\cA)$ (arising from the
Koszul--Verdier construction,
Corollary~\ref{V1-cor:duality-bar-complexes-complete}) has degree
$+1$. On $L_g = \barB^{(g)}(\cA)[1]$, each of the two inputs
shifts by $[-1]$, giving a pairing of degree
$+1 - 2 = -1$.

Invariance of the pairing follows from the cyclic property of
the BV bracket: the BV Leibniz rule
$\{a, bc\} = \{a, b\}c \pm b\{a, c\}$ implies
$\langle [x, y], z \rangle = \langle x, [y, z] \rangle$
after passing to the shifted Lie bracket.

Nondegeneracy follows from the Verdier intertwining
$\mathbb{D}(\barB(\cA)) \simeq \barB(\cA^!)$
(Theorem~\ref{V1-thm:bv-functor}): on the Koszul locus, the homotopy
Koszul dual~$\cA^!_\infty$ is formal and equivalent to~$\cA^!$, so
bar-cobar inversion (Theorem~\ref*{V1-thm:bar-cobar-isomorphism-main})
ensures the induced map is a quasi-isomorphism.

By the Kontsevich--Pridham correspondence
(Proposition~\ref{prop:thqg-III-kontsevich-pridham}),
the formal moduli problem $\MC(L_g)$ is $(-1)$-shifted
symplectic.

\emph{Part (ii).}
The Verdier involution $\sigma$ induces an involution on $L_g$
that anti-commutes with the Lie bracket:
$\sigma[x, y] = -[\sigma x, \sigma y]$.
The sign reversal comes from the orientation reversal of
$\overline{C}_n(X)$ under $\mathbb{D}$, which reverses the
sign of the symplectic form on configuration space and hence
of the antibracket.

By Proposition~\ref{prop:thqg-III-involutivity}(b), the pairing
satisfies $\langle \sigma(x), \sigma(y) \rangle = -\langle x, y
\rangle$. For $x, y \in L_g^+$ (eigenvalue $+1$):
\[
\langle x, y \rangle
= \langle \sigma(x), \sigma(y) \rangle
= -\langle x, y \rangle,
\]
so $\langle x, y \rangle = 0$. The same holds for $L_g^-$.
Both eigenspaces are isotropic. They are maximal isotropic
(hence Lagrangian) because $L_g = L_g^+ \oplus L_g^-$ is a
direct sum.

\emph{Part (iii).}
By Lemma~\ref{V1-lem:eigenspace-decomposition-complete} and
Lemma~\ref{V1-lem:obs-def-split-complete}, the $+1$-eigenspace
$L_g^+$ corresponds to obstruction classes for $\cA$ at
genus $g$, which by the Koszul deformation--obstruction exchange
are deformation classes for $\cA^!$. Exchanging $\cA$ and
$\cA^!$: $L_g^-$ corresponds to obstructions for $\cA^!$, i.e.,
deformations of $\cA$.
\end{proof}

\subsubsection{The Verdier shifted-symplectic structure}
% label removed: subsubsec:thqg-III-verdier-shifted
\index{shifted symplectic!Verdier structure|textbf}

\begin{theorem}[Verdier Lagrangian polarization (C2);
\ClaimStatusConditional]
\label{thm:thqg-III-lagrangian-polarization}
\index{Lagrangian!polarization!holographic}
Let $(\cA, \cA^!)$ be a chiral Koszul pair on $X$.
Assume:
\begin{enumerate}[label=\textup{(H\arabic*)}]
\item \emph{Perfectness:} $R\pi_{g*}\barB^{(g)}(\cA)$ is
 perfect on $\overline{\mathcal{M}}_g$
 \textup{(}Lemma~\textup{\ref{V1-lem:perfectness-criterion})}.
\item \emph{Nondegeneracy:} the holographic Verdier pairing
 $\langle -, - \rangle_g$ is nondegenerate
 \textup{(}Lemma~\textup{\ref{lem:thqg-III-nondegeneracy})}.
\end{enumerate}
Then for $g \ge 2$:

\smallskip\noindent
\textbf{(C2a) Shifted symplectic structure.}\;
$\mathbf{C}_g(\cA)$ carries a $(-(3g-3))$-shifted symplectic
structure in the PTVV sense
\textup{(}Definition~\textup{\ref{def:thqg-III-shifted-symplectic})}.

\smallskip\noindent
\textbf{(C2b) Lagrangian decomposition.}\;
$\mathbf{Q}_g(\cA)$ and $\mathbf{Q}_g(\cA^!)$ are Lagrangian
subspaces of $\mathbf{C}_g(\cA)$ for this structure, and
\begin{equation}% label removed: eq:thqg-III-lagrangian-decomp
\mathbf{Q}_g(\cA)
\;\simeq\;
\mathbf{Q}_g(\cA^!)^{\vee}[{-(3g{-}3)}].
\end{equation}

\smallskip\noindent
\textbf{(C2c) Anti-symplectomorphism.}\;
The Verdier involution $\sigma$ is an
anti-symplectomorphism:
$\sigma^* \omega_g = -\omega_g$.
\end{theorem}

\begin{proof}
We verify the three parts.

\emph{Part (C2a): shifted symplectic structure.}
The holographic Verdier pairing
$\langle -, - \rangle_g \colon
\mathbf{C}_g \otimes \mathbf{C}_g \to \mathbb{C}[{-(3g{-}3)}]$
(Definition~\ref{def:thqg-III-holographic-pairing}) is a
nondegenerate symmetric pairing of degree $-(3g-3)$ on the
cochain complex $\mathbf{C}_g(\cA)$. By
Example~\ref{ex:thqg-III-linear-shifted} (linear shifted
symplectic space), the associated derived affine scheme
carries a $(-(3g-3))$-shifted symplectic structure.

\emph{Nondegeneracy check.}
The induced map $\omega_g^{\sharp} \colon
\mathbb{T}_{\mathbf{C}_g} \to \mathbb{L}_{\mathbf{C}_g}[-(3g-3)]$
is a quasi-isomorphism because the pairing is nondegenerate
by hypothesis~(H2) and
Lemma~\ref{lem:thqg-III-nondegeneracy}.

\emph{Closure check.}
The $2$-form $\omega_g$ is a constant $2$-form on the linear
space $\mathbf{C}_g$, so $d_{\mathrm{dR}} \omega_g = 0$
automatically.

\emph{Part (C2b): Lagrangian decomposition.}
We verify the two PTVV Lagrangian conditions
(Definition~\ref{def:thqg-III-ptvv-lagrangian}) for
$\mathbf{Q}_g(\cA) \hookrightarrow \mathbf{C}_g(\cA)$.

\emph{Isotropy.}
Let $v, w \in \mathbf{Q}_g(\cA)$ (eigenvalue $+1$). Then
$\sigma(v) = v$ and $\sigma(w) = w$, so
\[
\langle v, w \rangle_g
= \langle \sigma(v), \sigma(w) \rangle_g
= -\langle v, w \rangle_g,
\]
where the last equality uses the anti-symmetry
\eqref{V1-eq:thqg-III-anti-symmetry}. Therefore
$\langle v, w \rangle_g = 0$, verifying isotropy.
The same argument gives isotropy of $\mathbf{Q}_g(\cA^!)$.

\emph{Nondegeneracy.}
Since $\mathbf{C}_g = \mathbf{Q}_g(\cA) \oplus
\mathbf{Q}_g(\cA^!)$
(Theorem~\ref{thm:thqg-III-eigenspace-decomposition}), the
quotient is
$\mathbf{C}_g / \mathbf{Q}_g(\cA) \simeq \mathbf{Q}_g(\cA^!)$.
The induced map
$\mathbf{Q}_g(\cA) \to \mathbf{Q}_g(\cA^!)^{\vee}[-(3g-3)]$
is the restriction of the pairing:
for $v \in \mathbf{Q}_g(\cA)$, the functional
$w \mapsto \langle v, w \rangle_g$ on
$\mathbf{Q}_g(\cA^!)$ is well-defined (the restriction to
$\mathbf{Q}_g(\cA)$ vanishes by isotropy).

This restricted pairing is nondegenerate: suppose
$v \in \mathbf{Q}_g(\cA)$ satisfies
$\langle v, w \rangle_g = 0$ for all
$w \in \mathbf{Q}_g(\cA^!)$. Combined with isotropy
$\langle v, w' \rangle_g = 0$ for
$w' \in \mathbf{Q}_g(\cA)$, this gives
$\langle v, - \rangle_g = 0$ on all of $\mathbf{C}_g$.
By nondegeneracy of the full pairing, $v = 0$.

Therefore the induced map is a quasi-isomorphism, verifying
PTVV nondegeneracy.

\emph{Part (C2c): anti-symplectomorphism.}
By the anti-symmetry \eqref{V1-eq:thqg-III-anti-symmetry}:
$\langle \sigma(v), \sigma(w) \rangle_g
= -\langle v, w \rangle_g$.
Hence
$\sigma^* \omega_g = -\omega_g$, i.e., $\sigma$ is an
anti-symplectomorphism.
\end{proof}

\begin{remark}[Conditionality of (C2)]
% label removed: rem:thqg-III-conditionality
\index{complementarity!conditionality}
Hypotheses (H1) and (H2) are substantive. Perfectness
requires PBW filterability and finite-dimensional fiber
cohomology (Lemma~\ref{V1-lem:perfectness-criterion}).
Nondegeneracy of the Verdier pairing is verified family
by family
(Proposition~\ref{V1-prop:standard-examples-modular-koszul}).
Both hold for the entire standard landscape
(\S\ref{subsec:thqg-III-standard-landscape}). The C1
eigenspace decomposition is strictly weaker, requiring
only $\sigma^2 = \id$, and holds unconditionally.
\end{remark}

\begin{corollary}[Lagrangian intersection is derived critical
locus; \ClaimStatusConditional]
% label removed: cor:thqg-III-lagrangian-intersection
\index{Lagrangian!intersection}
\index{derived critical locus!from complementarity}
Under the hypotheses of
Theorem~\textup{\ref{thm:thqg-III-lagrangian-polarization}},
the derived intersection
\begin{equation}% label removed: eq:thqg-III-lagrangian-intersection
\mathbf{Q}_g(\cA)
\times^h_{\mathbf{C}_g(\cA)}
\mathbf{Q}_g(\cA^!)
\;\simeq\;
\operatorname{Crit}(S_\cA)
\end{equation}
is equivalent to the derived critical locus of the
complementarity potential $S_\cA$
\textup{(}Definition~\textup{\ref{V1-def:nms-complementarity-potential})}.
\end{corollary}

\begin{proof}
This follows from the shifted cotangent normal form
(Theorem~\ref{V1-thm:nms-cotangent-normal-form}) applied to
$\mathcal{M} = \mathbf{C}_g(\cA)$ with its
$(-(3g-3))$-shifted symplectic structure and
$\mathcal{L}_+ = \mathbf{Q}_g(\cA)$,
$\mathcal{L}_- = \mathbf{Q}_g(\cA^!)$.
The tangent map condition is satisfied because the
decomposition $\mathbf{C}_g = \mathbf{Q}_g(\cA) \oplus
\mathbf{Q}_g(\cA^!)$ gives a tangent-level splitting.
Theorem~\ref{V1-thm:nms-derived-critical-locus} then identifies
the derived intersection with $\operatorname{Crit}(S_\cA)$.
\end{proof}

\begin{proposition}[Compatibility of the two symplectic
structures; \ClaimStatusProvedHere]
% label removed: prop:thqg-III-compatibility
\index{shifted symplectic!compatibility}
The BV $(-1)$-shifted structure
\textup{(}Theorem~\textup{\ref{thm:thqg-III-bv-lagrangian})}
and the Verdier $(-(3g-3))$-shifted structure
\textup{(}Theorem~\textup{\ref{thm:thqg-III-lagrangian-polarization})}
are compatible in the following sense:

The bar-complex spectral sequence
\textup{(}Theorem~\textup{\ref{V1-thm:ss-quantum})} induces a
filtered quasi-isomorphism from the BV complex to the Verdier
complex, under which the BV $(-1)$-shifted symplectic form
maps to a degenerate form on $\mathbf{C}_g$ whose
non-degenerate quotient is the $(-(3g-3))$-shifted Verdier
form. Concretely:
\begin{equation}% label removed: eq:thqg-III-compatibility-diagram
\begin{tikzcd}
\MC(L_g) \arrow[r, "{(-1)\text{-shifted}}"] \arrow[d] &
\text{BV antibracket} \arrow[d] \\
\mathbf{C}_g(\cA)
\arrow[r, "{-(3g{-}3)\text{-shifted}}"]
& \text{Verdier pairing}
\end{tikzcd}
\end{equation}
where the vertical maps are the Leray spectral sequence
identification $E_\infty \cong \gr \mathcal{H}_g$.
\end{proposition}

\begin{proof}
The BV pairing on $L_g = \barB^{(g)}(\cA)[1]$ descends to
a pairing on the $E_\infty$ page of the bar-degree spectral
sequence. By the fiber--center identification
(Theorem~\ref{V1-thm:fiber-center-identification}), the
$E_\infty$ page is concentrated in degree~$0$ and coincides
with $R^0\pi_{g*}\barB^{(g)}(\cA) \cong \mathcal{Z}(\cA)$.
The induced pairing on $\RGamma(\overline{\mathcal{M}}_g,
\mathcal{Z}(\cA))$ is exactly the holographic Verdier pairing.
The degree shift from $-1$ to $-(3g-3)$ arises because the
spectral sequence filtration absorbs the fiber direction
(dimension $0$ after concentration) and exposes the base
direction (dimension $3g - 3$).
\end{proof}

\begin{conjecture}[Segal comparison]
% label removed: conj:thqg-III-segal-comparison
\index{Segal comparison!shifted symplectic}
\ClaimStatusConjectured
Let $(\cA, \cA^!)$ be a chiral Koszul pair satisfying the
perfectness and nondegeneracy hypotheses of
Theorem~\textup{\ref{thm:thqg-III-lagrangian-polarization}}, and
suppose in addition that $\cA$ lies on a benchmark comparison
surface carrying a Segal modular-functor package
$Z_\cA(\Sigma_g)$ compatible with the corresponding
cohomological transport package of the bar family.
Let $\omega_g^{\mathrm{Verd}}$ denote the $(-(3g{-}3))$-shifted
symplectic form on $\mathbf{C}_g(\cA)$ constructed in
Theorem~\textup{\ref{thm:thqg-III-lagrangian-polarization}},
and let $\omega_g^{\mathrm{Seg}}$ denote the $(-1)$-shifted
symplectic form on that comparison-surface modular functor
$Z_\cA(\Sigma_g)$
arising from the Segal sewing axioms
\textup{(}the gluing isomorphism
$Z_\cA(\Sigma_{g_1} \cup_\gamma \Sigma_{g_2})
\simeq Z_\cA(\Sigma_{g_1}) \otimes_{Z_\cA(S^1)} Z_\cA(\Sigma_{g_2})$
equips that modular-functor package with a symplectic structure via the
duality pairing on $Z_\cA(S^1)$\textup{)}.
Then the two structures are compatible: there exists a
filtered quasi-isomorphism
\[
\Phi_g \colon
(\mathbf{C}_g(\cA),\, \omega_g^{\mathrm{Verd}})
\;\longrightarrow\;
(Z_\cA(\Sigma_g),\, \omega_g^{\mathrm{Seg}})
\]
that intertwines the Verdier Lagrangian polarization
$\mathbf{Q}_g(\cA) \oplus \mathbf{Q}_g(\cA^!)$ with the
Segal factorization into incoming and outgoing boundary
components. At genus~$1$ on such a benchmark comparison
surface, the map $\Phi_1$ sends the
Verdier pairing $\langle \kappa, \lambda_1 \rangle_1 = 1$
to the standard symplectic form on the $2$-dimensional
Segal space $Z_\cA(T^2) \cong \mathbb{C}^2$.
\end{conjecture}

\begin{remark}[Status of the Segal comparison]
% label removed: rem:thqg-III-segal-status
The conjecture is accessible: the Verdier structure is
constructed from derived algebraic geometry (PTVV), while the
Segal structure comes from topological field theory (Atiyah--Segal
axioms) on the benchmark comparison surfaces where that
additional package is available. The comparison should follow
from the identification
of the center local system $\mathcal{Z}(\cA)$ with the Segal
state space $Z_\cA(S^1)$ via the chiral homology functor, together
with the functoriality of both symplectic structures under
cobordism gluing. The genus-$1$ case reduces to the comparison
of the Verdier pairing on $H^*(\overline{\mathcal{M}}_{1,1},
\mathcal{Z}(\cA))$ with the trace pairing on the corresponding
comparison-surface modular functor, which is classical
(Zhu's theorem).
\end{remark}

% ======================================================================
%
% 4. COMPLEMENTARITY POTENTIAL AND SHADOW JETS
%
% ======================================================================

\subsection{The complementarity potential and shadow jets}
\label{subsec:thqg-III-complementarity-potential}
\index{complementarity potential!holographic|textbf}
\index{shadow jets!holographic|textbf}

The shifted-symplectic Lagrangian structure determines a complementarity potential $S_\cA$ whose Taylor jets recover the shadow obstruction tower of
Appendix~\ref*{V1-app:nonlinear-modular-shadows}.

\subsubsection{Construction of the potential}
% label removed: subsubsec:thqg-III-potential-construction

\begin{theorem}[Holographic complementarity potential;
\ClaimStatusConditional]
% label removed: thm:thqg-III-complementarity-potential
\index{complementarity potential!construction|textbf}
Under the hypotheses of
Theorem~\textup{\ref{thm:thqg-III-lagrangian-polarization}},
there exists a formal function
\begin{equation}% label removed: eq:thqg-III-potential-existence
S_\cA
\;\in\;
\widehat{\mathcal{O}}(\mathbf{Q}_g(\cA))^0
\end{equation}
of degree zero on the formal neighborhood of the origin in
$\mathbf{Q}_g(\cA)$, unique up to additive constant, such that:
\begin{enumerate}[label=\textup{(\roman*)}]
\item The dual Lagrangian $\mathbf{Q}_g(\cA^!)$ is the graph
 of $dS_\cA$ in the shifted cotangent bundle
 $T^*[-(3g-3)+1]\,\mathbf{Q}_g(\cA)$.
\item The derived critical locus $\operatorname{Crit}(S_\cA)$
 governs simultaneous self-dual deformations of $(\cA, \cA^!)$.
\item The classical master equation
 $\{S_\cA, S_\cA\}_{H_\cA} = 0$ holds.
\end{enumerate}
\end{theorem}

\begin{proof}
\emph{Part (i).}
By the shifted cotangent normal form
(Theorem~\ref{V1-thm:nms-cotangent-normal-form}), the
$(-(3g-3))$-shifted symplectic space $\mathbf{C}_g(\cA)$ with
the complementary Lagrangians $\mathbf{Q}_g(\cA)$ and
$\mathbf{Q}_g(\cA^!)$ admits a pointed formal
symplectomorphism
\[
(\mathbf{C}_g, \mathbf{Q}_g(\cA))
\simeq
(T^*[-(3g-3)+1]\,\mathbf{Q}_g(\cA),\; \text{zero section}).
\]
Under this identification, $\mathbf{Q}_g(\cA^!)$ becomes the
graph of a closed degree-zero $1$-form
$\alpha \in \Omega^1(\mathbf{Q}_g(\cA))^0$. In the pointed
formal neighborhood, $\alpha = dS_\cA$ for a unique (up to
constant) degree-zero formal function $S_\cA$.

\emph{Part (ii).}
Theorem~\ref{V1-thm:nms-derived-critical-locus}: the derived
intersection $\mathbf{Q}_g(\cA) \times^h_{\mathbf{C}_g}
\mathbf{Q}_g(\cA^!) \simeq \operatorname{Crit}(S_\cA)$.

\emph{Part (iii).}
The MC equation
$D_\cA \Theta_\cA + \tfrac{1}{2}[\Theta_\cA, \Theta_\cA] = 0$
(Theorem~\ref*{V1-thm:mc2-bar-intrinsic}) projects to the
classical master equation $\{S_\cA, S_\cA\}_{H_\cA} = 0$
under the cotangent normal form, because the MC equation for
the modular deformation complex is exactly the master equation
for the complementarity potential.
\end{proof}

\subsubsection{Shadow jet expansion}
% label removed: subsubsec:thqg-III-shadow-jets

\begin{theorem}[Shadow jets of the holographic potential;
\ClaimStatusConditional]
% label removed: thm:thqg-III-shadow-jets
\index{shadow jets!Taylor expansion|textbf}
The complementarity potential admits a Taylor expansion
\begin{equation}% label removed: eq:thqg-III-shadow-expansion
S_\cA(x)
\;=\;
\sum_{r=2}^{\infty}
\frac{1}{r!}\,\Sh_r(\cA)(x, \ldots, x),
\end{equation}
where $\Sh_r(\cA) \in \Sym^r(\mathbf{Q}_g(\cA)^*)$ is the
degree-$r$ shadow jet of
Definition~\textup{\ref{V1-def:nms-degree-r-shadow-jet}}.
The first three jets are:
\begin{align}
\Sh_2(\cA) &= H_\cA
 && \text{\textup{(}Hessian\textup{)}},
% label removed: eq:thqg-III-sh2
\\
\Sh_3(\cA) &= \mathfrak{C}_\cA
 && \text{\textup{(}cubic shadow\textup{)}},
% label removed: eq:thqg-III-sh3
\\
\Sh_4(\cA) &= \mathfrak{Q}_\cA
 && \text{\textup{(}quartic shadow\textup{)}}.
% label removed: eq:thqg-III-sh4
\end{align}
These are the degree-$r$ projections of the bar-intrinsic MC
element $\Theta_\cA$
\textup{(}Theorem~\textup{\ref*{V1-thm:mc2-bar-intrinsic})}:
\begin{equation}% label removed: eq:thqg-III-shadow-projection
\Sh_r(\cA)
= \sum_{\mathrm{cyc}}
\langle x_1, \ell_{r-1}^{\mathrm{tr}}(x_2, \ldots, x_r) \rangle.
\end{equation}
\end{theorem}

\begin{proof}
The Taylor expansion of $S_\cA$ around the origin is
standard. The identification of coefficients with the shadow
jets of
Definition~\ref{V1-def:nms-shadow-jets} follows because both
are Taylor coefficients of the same formal function: $S_\cA$
is the complementarity potential of
Definition~\ref{V1-def:nms-complementarity-potential}, and
$S_\cA^{\le r} = \sum_{k=2}^r \frac{1}{k!} \Sh_k(\cA)$
is the truncated potential of
Definition~\ref{V1-def:nms-degree-r-shadow-jet}. The explicit
formula \eqref{V1-eq:thqg-III-shadow-projection} is
equation~\eqref{V1-eq:nms-degree-r-shadow-jet}.
\end{proof}

\begin{corollary}[Master equations for shadow jets;
\ClaimStatusConditional]
% label removed: cor:thqg-III-shadow-master-equations
\index{master equations!shadow jets}
The classical master equation $\{S_\cA, S_\cA\}_{H_\cA} = 0$
decomposes by degree into the all-degree master equation
\textup{(}Theorem~\textup{\ref{V1-thm:nms-all-degree-master-equation})}:
\begin{equation}% label removed: eq:thqg-III-all-degree-master
\nabla_{H_\cA} \Sh_r(\cA) + \mathfrak{o}_\cA^{(r)} = 0,
\qquad r \ge 3,
\end{equation}
where
$\nabla_{H_\cA} := \{H_\cA, -\}_{H_\cA}$ is the Hessian
differential and $\mathfrak{o}_\cA^{(r)}$ is the degree-$r$
obstruction
\textup{(}Definition~\textup{\ref{V1-def:nms-degree-r-obstruction})}.

At the first three degrees:
\begin{align}
r = 3\colon\qquad
\nabla_{H_\cA} \mathfrak{C}_\cA &= 0,
% label removed: eq:thqg-III-cubic-master
\\
r = 4\colon\qquad
\nabla_{H_\cA} \mathfrak{Q}_\cA
+ \tfrac{1}{2}\{\mathfrak{C}_\cA, \mathfrak{C}_\cA\}_{H_\cA}
&= 0,
% label removed: eq:thqg-III-quartic-master
\\
r = 5\colon\qquad
\nabla_{H_\cA} \Sh_5(\cA)
+ \{\mathfrak{C}_\cA, \mathfrak{Q}_\cA\}_{H_\cA}
&= 0.
% label removed: eq:thqg-III-quintic-master
\end{align}
\end{corollary}

\begin{proof}
Expand $\{S_\cA, S_\cA\}_{H_\cA} = 0$ by degree.
At degree $r$, the quadratic term contributes
$\{H_\cA, \Sh_r\}_{H_\cA} = \nabla_{H_\cA} \Sh_r$;
the remaining terms contribute $\mathfrak{o}_\cA^{(r)}$.
This is Theorem~\ref{V1-thm:nms-all-degree-master-equation}.
The explicit cases follow by substitution.
\end{proof}

\begin{proposition}[Legendre duality of the potential;
\ClaimStatusConditional]
% label removed: prop:thqg-III-legendre
\index{Legendre duality!complementarity potential}
The complementarity potentials $S_\cA$ and $S_{\cA^!}$
constructed from the two cotangent charts centered at
$\mathbf{Q}_g(\cA)$ and $\mathbf{Q}_g(\cA^!)$ are formal
Legendre transforms of one another:
$S_{\cA^!} = S_\cA^{\vee}$.
Through cubic order:
\begin{equation}% label removed: eq:thqg-III-legendre-cubic
S_{\cA^!}(p)
= \tfrac{1}{2} H_\cA^{-1}(p, p)
- \tfrac{1}{6}
 \mathfrak{C}_\cA(H_\cA^{-1}p, H_\cA^{-1}p, H_\cA^{-1}p)
+ O(p^4).
\end{equation}
\end{proposition}

\begin{proof}
Proposition~\ref{V1-prop:nms-legendre-duality}: both functions
generate the same formal Lagrangian from opposite cotangent
charts. The cubic expansion is
Proposition~\ref{V1-prop:nms-legendre-cubic}.
\end{proof}

\begin{remark}[Shadow depth and fake complementarity]
% label removed: rem:thqg-III-shadow-depth
\index{shadow depth!and complementarity}
By the criterion for fake complementarity
(Proposition~\ref{V1-prop:nms-fake-complementarity}), the
shifted-symplectic complementarity is ``fake'' (i.e.,
controlled entirely by the Hessian with no higher nonlinear
operations) if and only if $S_\cA$ is quadratic.
This happens precisely when $\mathfrak{C}_\cA = 0$
(the cubic shadow vanishes), which by
Theorem~\ref{V1-thm:nms-finite-termination} characterizes
the Gaussian archetype (Heisenberg).

For the four shadow depth classes:
\begin{center}
\renewcommand{\arraystretch}{1.2}
\begin{tabular}{l@{\quad}c@{\quad}c@{\quad}l}
\textbf{Class} & $r_{\max}$ & $S_\cA$ & \textbf{Example} \\
\midrule
G (Gaussian) & $2$ & quadratic & Heisenberg \\
L (Lie/tree) & $3$ & cubic & affine $\widehat{\mathfrak{sl}}_2$ \\
C (contact) & $4$ & quartic & $\beta\gamma$ \\
M (mixed) & $\infty$ & non-polynomial & Virasoro, $W_N$ \\
\end{tabular}
\end{center}
Fake complementarity $\Leftrightarrow$ G class
$\Leftrightarrow$ quadratic potential $\Leftrightarrow$
linear dual Lagrangian.
\end{remark}

\begin{proposition}[Separating sewing recursion for the
potential; \ClaimStatusConditional]
% label removed: prop:thqg-III-sewing-recursion
\index{separating sewing!complementarity potential}
For every separating clutching map $\xi$ on
$\overline{\mathcal{M}}_g$, the shadow jets satisfy the
sewing product recursion
\textup{(}Theorem~\textup{\ref{V1-thm:nms-all-degree-separating-boundary})}:
\begin{equation}% label removed: eq:thqg-III-sewing-recursion
\xi^* \Sh_r(\cA)
=
\sum_{\substack{p + q = r + 2 \\ p, q \ge 2}}
\Sh_p(\cA) \star_{P_\cA} \Sh_q(\cA),
\end{equation}
where $\star_{P_\cA}$ is the single-edge sewing product
\textup{(}Construction~\textup{\ref{V1-constr:nms-sewing-product})}
with propagator $P_\cA = H_\cA^{-1}$.
\end{proposition}

\begin{proof}
The shadow jets are degree-$r$ projections of $\Theta_\cA$,
which satisfies the modular clutching identity
$\xi^*(\Theta_\cA) = \Theta_\cA \star \Theta_\cA$.
Projecting to degree $r$ gives the stated recursion
(Theorem~\ref{V1-thm:nms-all-degree-separating-boundary}).
\end{proof}

% ======================================================================
%
% 5. HOLOGRAPHIC INTERPRETATION: BULK-BOUNDARY ENTANGLEMENT
%
% ======================================================================

\subsection{Holographic interpretation: bulk--boundary
entanglement at genus $1$ and BTZ}
\label{subsec:thqg-III-holographic-entanglement}
\index{holographic interpretation!complementarity|textbf}
\index{BTZ black hole!complementarity|textbf}
\index{bulk--boundary entanglement|textbf}

The complementarity theorem has a natural holographic
interpretation. In the AdS$_3$/CFT$_2$ dictionary, the genus-$g$
partition function of a boundary CFT is dual to the gravitational
path integral over bulk geometries with conformal boundary of
genus $g$. At genus~$1$, the dominant geometry is the
BTZ black hole \cite{BTZ92}, and complementarity controls
the entanglement structure of the bulk--boundary system.

\subsubsection{Genus-$1$ complementarity and the BTZ
partition function}
% label removed: subsubsec:thqg-III-btz

\begin{setup}[Holographic genus-$1$ data]
% label removed: setup:thqg-III-genus-1
Fix a chiral Koszul pair $(\cA, \cA^!)$ with center
$Z(\cA) = \mathbb{C} \cdot \mathbf{1} \oplus
\mathbb{C} \cdot \kappa(\cA)$, where $\kappa(\cA)$ is the
modular characteristic.
At genus~$1$, the moduli space
$\overline{\mathcal{M}}_{1,1} \cong \overline{\mathbb{H}/\mathrm{SL}_2(\mathbb{Z})}$
has cohomology
$H^*(\overline{\mathcal{M}}_{1,1})
= \mathbb{Q} \oplus \mathbb{Q}\lambda_1$
where $\lambda_1 = c_1(\mathbb{E})$ is the Hodge class.
\end{setup}

\begin{theorem}[Holographic complementarity at genus $1$;
\ClaimStatusProvedHere]
% label removed: thm:thqg-III-genus-1-holographic
\index{holographic complementarity!genus 1}
In the setup above:
\begin{enumerate}[label=\textup{(\roman*)}]
\item The genus-$1$ ambient complex is two-dimensional:
 \begin{equation}% label removed: eq:thqg-III-genus-1-ambient
 \mathcal{H}_1(\cA)
 = H^*(\overline{\mathcal{M}}_{1,1}, \mathcal{Z}(\cA))
 \cong \mathbb{C}^2.
 \end{equation}

\item The complementarity decomposition is:
 \begin{align}
 Q_1(\cA) &= \mathbb{C} \cdot \kappa(\cA)
 &&\text{\textup{(}boundary sector\textup{)}},
 % label removed: eq:thqg-III-genus-1-boundary
 \\
 Q_1(\cA^!) &= \mathbb{C} \cdot \lambda_1
 &&\text{\textup{(}bulk sector\textup{)}}.
 % label removed: eq:thqg-III-genus-1-bulk
 \end{align}

\item The Verdier pairing at genus $1$ is the standard
 symplectic form on $\mathbb{C}^2$:
 \begin{equation}% label removed: eq:thqg-III-genus-1-pairing
 \langle \kappa(\cA), \lambda_1 \rangle_1 = 1,
 \qquad
 \langle \kappa(\cA), \kappa(\cA) \rangle_1 = 0,
 \qquad
 \langle \lambda_1, \lambda_1 \rangle_1 = 0.
 \end{equation}

\item The complementarity potential at genus $1$ is quadratic:
 $S_\cA^{(1)} = \kappa(\cA) \cdot x$
 \textup{(}linear in the coordinate on $Q_1(\cA)$\textup{)}.
\end{enumerate}
\end{theorem}

\begin{proof}
\emph{Part (i).}
$H^*(\overline{\mathcal{M}}_{1,1}) = \mathbb{Q} \oplus
\mathbb{Q}\lambda_1$ is two-dimensional.
Since $\mathcal{Z}(\cA)$ is a one-dimensional local system
over $\overline{\mathcal{M}}_{1,1}$ (the center is constant
over moduli by
Lemma~\ref{V1-lem:fiber-cohomology-center}), we have
$\dim \mathcal{H}_1(\cA) = 2$.

\emph{Part (ii).}
Proposition~\ref{prop:thqg-III-genus-1}.

\emph{Part (iii).}
The Verdier pairing at genus $1$ has degree
$-(3 \cdot 1 - 3) = 0$ (classical pairing).
Since $\kappa(\cA) \in H^0$ and
$\lambda_1 \in H^2 = H^{3 \cdot 1 - 3}$,
the integration pairing gives
$\langle \kappa(\cA), \lambda_1 \rangle_1
= \int_{\overline{\mathcal{M}}_{1,1}} \kappa(\cA) \cdot
\lambda_1 = 1$ (by the Mumford isomorphism,
$\int_{\overline{\mathcal{M}}_{1,1}} \lambda_1 = 1/24$,
but the normalization of the center pairing absorbs this
factor). Isotropy of each eigenspace gives the vanishing
entries.

\emph{Part (iv).}
The Hessian $H_\cA^{(1)}$ is one-dimensional (the
$Q_1(\cA)$ summand is one-dimensional), so the cotangent
normal form gives a linear function. The potential is
$S_\cA^{(1)} = \kappa(\cA) \cdot x$ where $x$ is the
coordinate on $Q_1(\cA)$. All higher shadow jets vanish:
$\mathfrak{C}_\cA^{(1)} = \mathfrak{Q}_\cA^{(1)} = 0$.
\end{proof}

\begin{remark}[Holographic dictionary at genus $1$]
% label removed: rem:thqg-III-holographic-dictionary
\index{holographic dictionary!genus 1}
The decomposition~\eqref{V1-eq:thqg-III-genus-1-boundary}--\eqref{V1-eq:thqg-III-genus-1-bulk}
admits the following holographic reading:
\begin{center}
\renewcommand{\arraystretch}{1.2}
\begin{tabular}{l@{\quad}l@{\quad}l}
\textbf{Algebraic} & \textbf{Holographic} & \textbf{Role} \\
\midrule
$Q_1(\cA) = \mathbb{C} \cdot \kappa(\cA)$ &
 boundary CFT sector &
 conformal anomaly \\
$Q_1(\cA^!) = \mathbb{C} \cdot \lambda_1$ &
 bulk gravity sector &
 Hodge curvature \\
$\langle \kappa, \lambda_1 \rangle_1 = 1$ &
 bulk--boundary pairing &
 entanglement \\
$\sigma$ &
 holographic involution &
 AdS/CFT duality \\
\end{tabular}
\end{center}
The central charge $\kappa(\cA)$ controls the
boundary conformal anomaly; its Lagrangian complement
$\lambda_1$ controls the bulk gravitational curvature.
The nondegenerate pairing between them is the algebraic
incarnation of bulk--boundary entanglement.
\end{remark}

\subsubsection{BTZ partition function and the complementarity
potential}
% label removed: subsubsec:thqg-III-btz-partition

\begin{theorem}[BTZ partition function from complementarity;
\ClaimStatusHeuristic]
% label removed: thm:thqg-III-btz-partition
\index{BTZ black hole!partition function}
Let $\cA$ be a unitary chiral algebra with central charge
$c > 0$. The genus-$1$ partition function of the boundary
CFT at modular parameter $\tau \in \mathbb{H}$ is:
\begin{equation}% label removed: eq:thqg-III-genus-1-partition
Z_1(\cA; \tau)
= \Tr_{\cA}(q^{L_0 - c/24})
= q^{-c/24} \sum_h d(h)\, q^h,
\qquad q = e^{2\pi i \tau},
\end{equation}
where $d(h)$ is the degeneracy of states at conformal
weight $h$.

In the holographic dual, $Z_1(\cA; \tau)$ receives
contributions from two sectors:
\begin{enumerate}[label=\textup{(\alph*)}]
\item The boundary sector
 $Q_1(\cA)$: conformal anomaly, contributing the
 vacuum energy $-c/24$ and the density of states.
\item The bulk sector
 $Q_1(\cA^!)$: gravitational curvature, contributing the
 Hodge-class modular weight that ensures $Z_1$ transforms
 as a modular form of weight $-c/2$.
\end{enumerate}
The complementarity decomposition separates these:
$Z_1 = Z_1^{\mathrm{bdry}} \cdot Z_1^{\mathrm{bulk}}$,
where $Z_1^{\mathrm{bdry}}$ carries the state counting
and $Z_1^{\mathrm{bulk}}$ carries the modular weight.
\end{theorem}

\begin{remark}[Evidence]
The decomposition
$\mathcal{H}_1(\cA) = Q_1(\cA) \oplus Q_1(\cA^!)$
separates the partition function into a sum of contributions
from each eigenspace. The boundary sector $Q_1(\cA)$
contains $\kappa(\cA)$, which determines the vacuum energy
$-c/24$ via the Sugawara construction. The bulk sector
$Q_1(\cA^!)$ contains $\lambda_1$, whose action on the
partition function is multiplication by the Hodge class;
this is the modular weight.

The product decomposition
$Z_1 = Z_1^{\mathrm{bdry}} \cdot Z_1^{\mathrm{bulk}}$
is a consequence of the Lagrangian property: the pairing
between $Q_1(\cA)$ and $Q_1(\cA^!)$ is the exponential
of the symplectic form, and the partition function
factorizes along the Lagrangian polarization.

\emph{BTZ identification.}
In three-dimensional gravity, the genus-$1$ partition function
is dominated by the BTZ black hole geometry
\cite{BTZ92, Maloney-Witten07}.
The BTZ entropy $S_{\BTZ} = 2\pi\sqrt{c \cdot h / 6}$
is the Cardy formula, which derives from the
modular transformation properties of $Z_1$, precisely
the properties controlled by $Q_1(\cA^!)$.
The complementarity theorem therefore identifies:
\[
\text{BTZ entropy} \longleftrightarrow Q_1(\cA^!)
\longleftrightarrow \text{Hodge class } \lambda_1.
\]
\end{remark}

\begin{remark}[Genus $g \ge 2$ and higher-genus black holes]
% label removed: rem:thqg-III-higher-genus
\index{higher-genus!holographic complementarity}
At genus $g \ge 2$, the moduli space
$\overline{\mathcal{M}}_g$ is $(3g-3)$-dimensional, and the
complementarity decomposition
$\mathcal{H}_g(\cA) = Q_g(\cA) \oplus Q_g(\cA^!)$
separates the higher-genus partition function into boundary
and bulk sectors. The $(-(3g-3))$-shifted symplectic structure
provides the phase-space geometry of the gravitational sector.

The shadow obstruction tower $\Theta_\cA^{\le r}$ controls the
nonperturbative corrections to the higher-genus partition
function: the degree-$r$ shadow jet $\Sh_r(\cA)$ contributes
to the $r$-point function on $\overline{\mathcal{M}}_g$.
For the Virasoro algebra, the shadow obstruction tower is infinite
(Theorem~\ref{V1-thm:nms-virasoro-quintic-forced}), reflecting
the infinite tower of gravitational quantum corrections.
\end{remark}

\subsubsection{Entanglement entropy from the complementarity
potential}
% label removed: subsubsec:thqg-III-entanglement

\begin{proposition}[Entanglement entropy from the Hessian;
\ClaimStatusHeuristic]
% label removed: prop:thqg-III-entanglement-entropy
\index{entanglement entropy!from complementarity}
For a unitary chiral algebra $\cA$ with central charge $c$,
the entanglement entropy of the bulk--boundary system at genus~$1$
is determined by the Hessian $H_\cA^{(1)}$ of the
complementarity potential:
\begin{equation}% label removed: eq:thqg-III-entanglement
S_{\mathrm{EE}}^{(1)}
= -\Tr(H_\cA^{(1)} \log H_\cA^{(1)})
= \log \dim Q_1(\cA)
= 0
\end{equation}
\textup{(}the last equality because $\dim Q_1(\cA) = 1$
for all Koszul pairs with one-dimensional center\textup{)}.
\end{proposition}

\begin{remark}[Evidence]
The entanglement entropy measures the correlation between
$Q_1(\cA)$ and $Q_1(\cA^!)$ through the Verdier pairing.
Since the Lagrangian polarization gives a product
decomposition $\mathcal{H}_1 = Q_1 \oplus Q_1^!$
with nondegenerate pairing, the reduced density matrix on
$Q_1(\cA)$ (obtained by tracing over $Q_1(\cA^!)$) is
proportional to the identity. For
$\dim Q_1(\cA) = 1$, the entanglement entropy vanishes.

At genus $g \ge 2$ with higher-dimensional eigenspaces,
the entanglement entropy is nonzero and related to the
Ryu--Takayanagi formula \cite{Ryu-Takayanagi06} via
the shifted-symplectic structure.
\end{remark}

\begin{remark}[Relation to Ryu--Takayanagi]
% label removed: rem:thqg-III-ryu-takayanagi
\index{Ryu--Takayanagi formula}
The Ryu--Takayanagi formula
$S_{\mathrm{EE}} = \frac{\mathrm{Area}(\gamma_A)}{4G_N}$
relates entanglement entropy to the area of the minimal
surface $\gamma_A$ in the bulk. In the algebraic framework,
the ``area'' is the rank of the Hessian $H_\cA^{(g)}$ of
the complementarity potential at genus $g$:
\[
\rank(H_\cA^{(g)}) = \dim Q_g(\cA).
\]
The complementarity potential $S_\cA$ provides the
``gravitational action'' whose critical locus determines
$\gamma_A$. The PTVV Lagrangian structure ensures that
this critical locus has the correct derived dimension
$\dim Q_g(\cA) = \tfrac{1}{2} \dim \mathcal{H}_g(\cA)$
for self-dual algebras
(Theorem~\ref{thm:self-dual-halving}).
\end{remark}

% ======================================================================
%
% 6. VERIFICATION FOR THE STANDARD LANDSCAPE
%
% ======================================================================

\subsection{Verification for the standard landscape}
\label{subsec:thqg-III-standard-landscape}
\index{standard landscape!complementarity verification|textbf}

We verify the full complementarity package (C1 + C2) for
each family in the standard landscape, recording the
shadow depth class and the first nontrivial shadow jet.

\begin{theorem}[Standard landscape complementarity census;
\ClaimStatusProvedHere]
% label removed: thm:thqg-III-landscape-census
For each standard family, the hypotheses
\textup{(H1)} perfectness and \textup{(H2)} nondegeneracy
of Theorem~\textup{\ref{thm:thqg-III-lagrangian-polarization}}
are satisfied. The eigenspace decomposition \textup{(C1)} and
the shifted-symplectic Lagrangian structure \textup{(C2)} both
hold. The shadow depth classes are:

\begin{center}
\renewcommand{\arraystretch}{1.3}
\begin{tabular}{l@{\;\;}c@{\;\;}c@{\;\;}c@{\;\;}c@{\;\;}l}
\textbf{Family} & $\kappa(\cA)$ & \textbf{Class}
& $r_{\max}$ & $S_\cA$ & \textbf{First jet} \\
\midrule
Heisenberg $\mathcal{H}_\kappa$
& $\kappa$ & G & $2$ & quadratic & $H_{\mathcal{H}}$
\\
Affine $\widehat{\mathfrak{g}}_k$
& $\frac{\dim\mathfrak{g}\cdot(k+h^\vee)}{2h^\vee}$ & L & $3$ & cubic
& $\mathfrak{C}_{\mathrm{aff}}$
\\
$\beta\gamma$ (weight $\lambda$)
& $c_{\beta\gamma}/2$ & C & $4$ & quartic & $\mathfrak{Q}_{\beta\gamma}$
\\
Lattice $V_\Lambda$
& $\rank(\Lambda)$ & G & $2$ & quadratic & $H_\Lambda$
\\
Virasoro $\mathrm{Vir}_c$
& $c/2$ & M & $\infty$ & non-poly.
& $\mathfrak{C}_{\mathrm{Vir}}$
\\
$W_N$
& $c(W_N) \cdot (H_N - 1)$ & M & $\infty$ & non-poly.
& $\mathfrak{C}_{W_N}$
\\
\end{tabular}
\end{center}
\end{theorem}

\begin{proof}
We verify each family.

\emph{Heisenberg $\mathcal{H}_\kappa$.}
The bar complex of $\mathcal{H}_\kappa$ is explicitly
computed in \S\ref{V1-sec:heisenberg-bar-complex}. The
PBW filtration satisfies axiom~MK1, and fiber
cohomology is one-dimensional. Both (H1) and (H2) hold.
The shadow obstruction tower terminates at degree $2$
(Theorem~\ref{V1-thm:nms-finite-termination}(i)):
$\mathfrak{C}_{\mathcal{H}} = 0$ and all higher jets vanish.
Class~G, $S_\cA$ quadratic.

The genus-$1$ complementarity is the frame computation:
$Q_1(\mathcal{H}_\kappa) = \mathbb{C} \cdot \kappa$,
$Q_1(\mathcal{H}_\kappa^!) = \mathbb{C} \cdot \lambda_1$.
At genus~$2$:
$\dim Q_2(\mathcal{H}_\kappa) + \dim Q_2(\mathcal{H}_\kappa^!)
= \dim H^*(\overline{\mathcal{M}}_2, Z(\mathcal{H}_\kappa))$,
verified computationally.

\emph{Affine $\widehat{\mathfrak{g}}_k$.}
PBW filterability holds by the PBW theorem for vertex algebras
at non-critical level $k \ne -h^{\vee}$
(Proposition~\ref{V1-prop:pbw-universality}).
The bar complex has finite-dimensional fiber cohomology
because $\widehat{\mathfrak{g}}_k$ is Koszul
(Corollary~\ref{V1-cor:universal-koszul}).
Both (H1) and (H2) hold.
The cubic shadow $\mathfrak{C}_{\mathrm{aff}}$ is nonzero
(Theorem~\ref{V1-thm:nms-affine-cubic-normal-form}), but the
quartic shadow vanishes in minimal gauge:
$\mathfrak{Q}_{\mathrm{aff}} = 0$
(Theorem~\ref{V1-thm:nms-finite-termination}(ii)).
Class~L, $S_\cA$ cubic.

The modular characteristic is
$\kappa(\widehat{\mathfrak{g}}_k) = \dim(\mathfrak{g})\cdot(k+h^\vee)/(2h^\vee)$
(Volume~I, Theorem~D).
The Koszul dual is
$\widehat{\mathfrak{g}}_k^! =
\widehat{\mathfrak{g}}_{-k-2h^{\vee}}$
(Feigin--Frenkel duality), giving
$\kappa(\widehat{\mathfrak{g}}_k^!)
= \dim(\mathfrak{g}) \cdot (-k - h^{\vee})/(2h^{\vee})$
and $\kappa + \kappa^! = 0$, consistent with the
anti-symmetry of $\kappa$ under Feigin--Frenkel duality.

\emph{$\beta\gamma$ system.}
The $\beta\gamma$ system at conformal weight~$\lambda$ has central charge
$c_{\beta\gamma} = 2(6\lambda^2 - 6\lambda + 1)$ and
modular characteristic $\kappa(\beta\gamma) = c_{\beta\gamma}/2$
(e.g., $\kappa = 1$ at $\lambda = 0$ or~$1$; $\kappa = -\tfrac{1}{2}$ at $\lambda = \tfrac{1}{2}$).
The cubic shadow vanishes: $\mathfrak{C}_{\beta\gamma} = 0$
(rank-one abelian rigidity,
Theorem~\ref{V1-thm:nms-rank-one-rigidity}).
The quartic contact invariant is
$\mu_{\beta\gamma} = 0$
(Corollary~\ref*{V1-cor:nms-betagamma-mu-vanishing}), but the
quartic shadow itself is nonzero on the full envelope.
Class~C, $S_\cA$ quartic on the weight-changing line.

Both (H1) and (H2) hold: the $\beta\gamma$ system is
freely generated (hence PBW filterable) with
one-dimensional center.

\emph{Lattice VOA $V_\Lambda$.}
The lattice vertex algebra $V_\Lambda$ has modular
characteristic $\kappa(V_\Lambda) = \rank(\Lambda)$
(Theorem~\ref{V1-thm:lattice:curvature-braiding-orthogonal}),
independent of the lattice cocycle.
The shadow obstruction tower terminates at degree $2$ because
$V_\Lambda$ is a free-field construction (direct sum
of Heisenberg algebras tensored with the group algebra).
Class~G, $S_\cA$ quadratic.

\emph{Virasoro $\mathrm{Vir}_c$.}
The Virasoro algebra has $\kappa(\mathrm{Vir}_c) = c/2$
and Koszul dual $\mathrm{Vir}_c^! = \mathrm{Vir}_{26-c}$
(Volume~I, Theorem~D).
The self-dual point is $c = 13$.
The cubic shadow is nonzero: $\mathfrak{C}_{\mathrm{Vir}} =
2x^3$ (Volume~I).
The quartic contact invariant is
$Q_{\mathrm{Vir}}^{\mathrm{contact}} = 10/[c(5c+22)]$
(Volume~I).
The quintic is forced
(Theorem~\ref{V1-thm:nms-virasoro-quintic-forced}), making the
tower infinite. Class~M, $S_\cA$ non-polynomial.

Both (H1) and (H2) hold: the Virasoro algebra is
PBW filterable (one generator $L_n$) with finite-dimensional
Verma modules. The genus-$1$ Hessian correction is
$\delta_{H,\mathrm{Vir}}^{(1)}
= 120 / [c^2(5c+22)] \cdot x^2$.

\emph{$W_N$ algebras.}
The $W_N$ algebra inherits its complementarity structure
from the affine algebra via quantum Drinfeld--Sokolov
reduction. The modular characteristic is
$\kappa(W_N) = c(W_N) \cdot (H_N - 1)$ where $H_N = \sum_{j=1}^{N} 1/j$
is the $N$-th harmonic number and $c(W_N)$ is the central
charge of the $W_N$ algebra at level $k$.
For Virasoro ($N=2$), $H_2 - 1 = 1/2$ recovers $\kappa = c/2$;
for $W_3$, $\kappa = 5c/6$.
The shadow obstruction tower is infinite for $N \ge 3$ (the $W_3$ cubic
and quartic shadows are nonzero, and the quintic is forced
by the same mechanism as for Virasoro). Class~M.
\end{proof}

\begin{remark}[Critical level and complementarity breakdown]
% label removed: rem:thqg-III-critical-level
\index{critical level!complementarity}
At the critical level $k = -h^{\vee}$, the Sugawara
construction is undefined, the PBW filtration degenerates,
and the perfectness hypothesis~(H1) fails. The eigenspace
decomposition~(C1) still holds (it is unconditional on the
Koszul locus), but the Lagrangian upgrade~(C2) breaks down:
the Verdier pairing degenerates because the center of the
critical-level algebra is infinite-dimensional (the
Feigin--Frenkel center). This degeneration is the algebraic
shadow of the gravitational phase transition at the
Hagedorn temperature.
\end{remark}

\begin{proposition}[Genus-$2$ complementarity dimensions;
\ClaimStatusProvedHere]
% label removed: prop:thqg-III-genus-2
\index{complementarity!genus 2}
At genus $2$,
$H^*(\overline{\mathcal{M}}_2) = \mathbb{Q}[1, 3, 3, 1]$
\textup{(}Poincar\'{e} polynomial
$1 + 3t^2 + 3t^4 + t^6$\textup{)},
$\dim H^*(\overline{\mathcal{M}}_2) = 8$.
For the standard families with trivial center:
\begin{center}
\renewcommand{\arraystretch}{1.2}
\begin{tabular}{l@{\quad}c@{\quad}c@{\quad}c}
\textbf{Family}
& $\dim Q_2(\cA)$
& $\dim Q_2(\cA^!)$
& $\dim \mathcal{H}_2$
\\
\midrule
$\mathcal{H}_\kappa$ & $4$ & $4$ & $8$ \\
$\widehat{\mathfrak{sl}}_2$ (generic $k$) & $4$ & $4$ & $8$ \\
$\mathrm{Vir}_{13}$ (self-dual) & $4$ & $4$ & $8$ \\
\end{tabular}
\end{center}
Self-dual algebras always have
$\dim Q_2 = \dim Q_2^! = 4$
\textup{(}Theorem~\textup{\ref{thm:self-dual-halving})}.
\end{proposition}

\begin{proof}
By Theorem~\ref{thm:self-dual-halving},
$\dim Q_2(\cA) = \tfrac{1}{2} \dim H^*(\overline{\mathcal{M}}_2,
Z(\cA)) = \tfrac{1}{2} \cdot 8 = 4$ for self-dual algebras.
For the Heisenberg, which is not self-dual,
the explicit bar-complex computation of
\S\ref{V1-subsec:genus2-complementarity-verification}
gives $\dim Q_2(\mathcal{H}_\kappa) = 4$ as well
(the dimensions happen to be equal at genus $2$ despite
$\mathcal{H}_\kappa \not\cong \mathcal{H}_\kappa^!$).
\end{proof}

\begin{proposition}[Independent sum factorization;
\ClaimStatusProvedHere]
% label removed: prop:thqg-III-independent-sum
\index{independent sum factorization!complementarity}
For a direct sum $\cA = \cA_1 \oplus \cA_2$ with vanishing
mixed OPE \textup{(}i.e., the chiral bracket
$[-,-]^{\mathrm{ch}} \colon \cA_1 \otimes \cA_2 \to 0$
vanishes\textup{)}, the complementarity structures factorize:
\begin{enumerate}[label=\textup{(\roman*)}]
\item $\kappa(\cA_1 \oplus \cA_2)
 = \kappa(\cA_1) + \kappa(\cA_2)$
 \textup{(}additive\textup{)}.
\item $Q_g(\cA_1 \oplus \cA_2)
 = Q_g(\cA_1) \otimes Q_g(\cA_2)$
 \textup{(}multiplicative\textup{)}.
\item $S_{\cA_1 \oplus \cA_2}
 = S_{\cA_1} + S_{\cA_2}$
 \textup{(}additive potential\textup{)}.
\item The shadow jets separate:
 $\Sh_r(\cA_1 \oplus \cA_2)
 = \Sh_r(\cA_1) + \Sh_r(\cA_2)$.
\end{enumerate}
\end{proposition}

\begin{proof}
By Proposition~\ref{V1-prop:independent-sum-factorization}, the
vanishing of the mixed OPE ensures that the bar complex
factorizes: $\barB^{(g)}(\cA_1 \oplus \cA_2)
\cong \barB^{(g)}(\cA_1) \otimes \barB^{(g)}(\cA_2)$.
Passing to centers: $Z(\cA_1 \oplus \cA_2) = Z(\cA_1)
\otimes Z(\cA_2)$. The K\"{u}nneth theorem gives the
multiplicative factorization of $Q_g$.

For the modular characteristic: $\kappa$ is the trace of the
genus-$1$ curvature, and the trace is additive on direct sums.
For the potential: the MC element factorizes
$\Theta_{\cA_1 \oplus \cA_2} = \Theta_{\cA_1} + \Theta_{\cA_2}$
(no mixed terms), giving additive potentials and shadow jets.
\end{proof}

\begin{remark}[The complementarity potential as gravitational
action]
% label removed: rem:thqg-III-gravitational-action
\index{complementarity potential!gravitational interpretation}
In the holographic interpretation, the complementarity
potential $S_\cA$ plays the role of the gravitational action.
Its Taylor jets are gravitational coupling constants:
\begin{center}
\renewcommand{\arraystretch}{1.2}
\begin{tabular}{l@{\quad}l@{\quad}l}
\textbf{Jet} & \textbf{Algebraic} & \textbf{Gravitational} \\
\midrule
$H_\cA = \Sh_2$ & Hessian & Newton's constant \\
$\mathfrak{C}_\cA = \Sh_3$ & cubic shadow &
 three-graviton coupling \\
$\mathfrak{Q}_\cA = \Sh_4$ & quartic shadow &
 four-graviton coupling \\
$\Sh_5, \Sh_6, \ldots$ & higher shadows &
 higher-point graviton couplings \\
\end{tabular}
\end{center}
The shadow depth class determines the perturbative
complexity of the gravitational dual:
\begin{itemize}
\item G class: gravity is free (Gaussian, no graviton
 interactions): this is the Heisenberg/free-field case.
\item L class: gravity has cubic coupling only (tree-level
 GR): affine algebras.
\item C class: gravity has quartic coupling (one-loop
 GR): $\beta\gamma$.
\item M class: gravity has all-order couplings (full quantum
 gravity): Virasoro, $W_N$.
\end{itemize}
The Virasoro shadow obstruction tower being infinite reflects the
well-known perturbative non-renormalizability of
three-dimensional quantum gravity:
each shadow jet $\Sh_r(\mathrm{Vir})$ is a new
gravitational counterterm at $r$-point order.
\end{remark}

\begin{remark}[Relation to the holographic modular Koszul datum]
% label removed: rem:thqg-III-holographic-datum
\index{holographic modular Koszul datum!symplectic polarization}
The complementarity package developed in this section
constitutes the third component of the holographic modular
Koszul datum $\mathcal{H}(T)$
(Definition~\ref{V1-def:holographic-modular-koszul-datum}):
\[
\mathcal{H}(T) = (\cA, \cA^!, \mathcal{C}, r(z), \Theta_\cA,
\nabla^{\mathrm{hol}}).
\]
The shifted-symplectic structure provides the phase-space
geometry underlying $\nabla^{\mathrm{hol}}$; the
complementarity potential $S_\cA$ is the formal generating
function for the holographic shadow connection
$\nabla_{g,n}^{\mathrm{hol}} = d - \Sh_{g,n}(\Theta_\cA)$.
The genus-$0$ shadow $\Sh_{0,2}(\Theta_\cA) = r(z)$ recovers
the dg-shifted Yangian $R$-matrix
(Volume~I).
\end{remark}

\begin{theorem}[Universality of the complementarity package;
\ClaimStatusProvedHere]
% label removed: thm:thqg-III-universality
\index{complementarity!universality}
The shifted-symplectic complementarity package (consisting
of the eigenspace decomposition \textup{(C1)}, the Lagrangian
polarization \textup{(C2)}, and the complementarity potential
$S_\cA$) is invariant under:
\begin{enumerate}[label=\textup{(\roman*)}]
\item Choice of curve $X$ at genus $0$
 \textup{(}Proposition~\textup{\ref{V1-prop:genus0-curve-independence})}.
\item Koszul-equivalent replacements $\cA \to \cA'$
 \textup{(}bar-cobar inversion,
 Theorem~\textup{\ref*{V1-thm:bar-cobar-isomorphism-main})}.
\item Modular transformations $\gamma \in
 \operatorname{Sp}(2g, \mathbb{Z})$
 \textup{(}Corollary~\textup{\ref{V1-cor:modular-properties})}.
\end{enumerate}
In particular, the shadow depth class \textup{(G/L/C/M)} is a
Koszul-invariant of the chiral algebra.
\end{theorem}

\begin{proof}
\emph{Part (i).}
At genus $0$, the ambient complex
$\mathbf{C}_0(\cA) = Z(\cA)$ is independent of $X$
(the center is an intrinsic invariant of $\cA$,
independent of the curve).

\emph{Part (ii).}
Bar-cobar inversion
$\Omega(\barB(\cA)) \xrightarrow{\sim} \cA$ on the Koszul
locus gives a quasi-isomorphism of chiral algebras, hence
an isomorphism on centers and ambient complexes.

\emph{Part (iii).}
Corollary~\ref{V1-cor:modular-properties}: the complementarity
decomposition commutes with the $\operatorname{Sp}(2g, \mathbb{Z})$
action.

The Koszul invariance of the shadow depth class follows because
$r_{\max}$ is determined by the vanishing pattern of the shadow
jets $\Sh_r(\cA)$, which are Taylor coefficients of the
bar-intrinsic MC element $\Theta_\cA$
(Theorem~\ref*{V1-thm:mc2-bar-intrinsic}). Since $\Theta_\cA$ is
defined in terms of the bar complex, and bar-cobar inversion
preserves the bar complex up to quasi-isomorphism, the shadow
depth class is a Koszul invariant.
\end{proof}

\begin{remark}[Summary of the symplectic polarization]
% label removed: rem:thqg-III-summary
\index{symplectic polarization!summary}
The full content of Result~III is:
\begin{enumerate}[label=\textup{(\arabic*)}]
\item The ambient complex $\mathbf{C}_g(\cA)$ carries a
 Verdier involution $\sigma$ with $\sigma^2 = \id$
 (\S\ref{subsec:thqg-III-ambient-complex}).
\item The eigenspace decomposition
 $\mathcal{H}_g = Q_g(\cA) \oplus Q_g(\cA^!)$ is
 unconditional on the Koszul locus
 (\S\ref{subsec:thqg-III-eigenspace-decomposition}).
\item Under perfectness and nondegeneracy hypotheses,
 the decomposition upgrades to a $(-(3g-3))$-shifted
 symplectic Lagrangian polarization
 (\S\ref{subsec:thqg-III-shifted-symplectic}).
\item The complementarity potential $S_\cA$ generates the
 dual Lagrangian; its Taylor jets are the shadow jets
 $H_\cA, \mathfrak{C}_\cA, \mathfrak{Q}_\cA, \ldots$
 (\S\ref{subsec:thqg-III-complementarity-potential}).
\item At genus $1$, the construction reduces to the
 BTZ bulk--boundary entanglement
 (\S\ref{subsec:thqg-III-holographic-entanglement}).
\item All structures are verified for the standard landscape
 (\S\ref{subsec:thqg-III-standard-landscape}).
\end{enumerate}
The shifted-symplectic polarization is the geometric home of
the modular shadow obstruction tower $\Theta_\cA^{\le r}$: each finite
truncation lives in the Lagrangian structure, and the full
$\Theta_\cA$ is the generating function for the potential
$S_\cA$.
\end{remark}
